\documentclass[11pt,reqno]{amsart}

\usepackage[margin=1in]{geometry}
\usepackage[T1]{fontenc}
\usepackage{lmodern}
\usepackage{microtype}
\usepackage{mathtools}
\usepackage{amssymb}
\usepackage{xcolor}
\usepackage{braket}
\usepackage[
  colorlinks=true,
  linkcolor=blue!55!black,
  citecolor=blue!55!black,
  urlcolor=blue!55!black
]{hyperref}
\allowdisplaybreaks

\newtheorem{theorem}{Theorem}[section]
\newtheorem{corollary}[theorem]{Corollary}
\newtheorem{proposition}[theorem]{Proposition}
\newtheorem{lemma}[theorem]{Lemma}
\theoremstyle{remark}
\newtheorem{remark}[theorem]{Remark}

\usepackage[nameinlink,capitalize,noabbrev]{cleveref}

\DeclareMathOperator{\Tr}{Tr}
\DeclareMathOperator{\rank}{rank}
\DeclareMathOperator{\SN}{SN}
\DeclareMathOperator{\SR}{SR}
\DeclareMathOperator{\ran}{ran}

\newcommand{\cL}{\mathcal L}
\newcommand{\id}{\mathrm{id}}
\newcommand{\ketbra}[2]{\lvert #1\rangle\!\langle #2\rvert}
\newcommand{\ip}[2]{\langle #1,#2\rangle}
\newcommand{\norm}[1]{\lVert #1\rVert}
\newcommand{\abs}[1]{\lvert #1\rvert}
\newcommand{\C}{\mathbb C}

\title[A Sharp Rank-$r$ Partial-Trace Inequality]
{A Sharp Rank-$r$ Partial-Trace Inequality via Double-Skew Singular-Value Bounds}
\author{Elazar Gershuni}
\date{July 26, 2026}

\hypersetup{
  pdftitle={A Sharp Rank-r Partial-Trace Inequality via Double-Skew Singular-Value Bounds},
  pdfauthor={Elazar Gershuni},
  pdfsubject={Rank-constrained partial-trace inequalities and applications}
}

\begin{document}

\begin{abstract}
Using the recent sharp Schmidt-rank-two estimate of Fu, Gao,
and Park for a double-antisymmetric projection, we prove the
arbitrary-matrix rank-$r$ partial-trace inequality for all ranks. Specifically, for every finite-dimensional bipartite Hilbert
space and every operator $C$ of rank at most $r$,
\[
    \norm{\Tr_U C}_2^2+\norm{\Tr_V C}_2^2
    \le r\norm{C}_2^2+\frac{1}{r}\abs{\Tr C}^2.
\]
The proof constructs an exact completely positive map decomposition for
the negative partial-transpose sector and controls the resulting terms
through a complete-graph edge summation over
$\mathfrak{so}(U)\otimes\mathfrak{so}(V)$. We derive the exact
rank-constrained two-copy score curve and asymmetric tensor-product
threshold, convert the critical score into a four-qudit one-sided
Schmidt-number witness with perfect completeness, and formulate the
higher-copy problem through an exact code-state operator. We also
record the complete-graph Gram defect underlying the proof, together
with selected consequences for positive maps, local witnesses, and
Kronecker-sum singular values.
\end{abstract}

\maketitle

\begin{quote}\small
\textbf{AI-generation and verification disclosure.}
The problem selection, prompting, and iterative direction were supplied
by the author. The mathematical arguments, claimed results,
applications, and exposition were generated by OpenAI's GPT-5.6.
The author did not independently derive the proofs, and no claim of
independent human verification is made. The manuscript has not received independent expert review and should not
be relied upon until its arguments have been checked by qualified experts.
\end{quote}

\section{Introduction}

Costa Rico proved the rank-sensitive partial-trace inequality below
for normal operators and related its unrestricted rank-two case to
two-copy distillability of Werner states \cite{CostaRico}. Costa Rico
and Wolf subsequently developed partial-trace relations for general
matrices, including singular-value majorization results for Kronecker
sums and applications to Schmidt-number witnesses and positive maps,
while identifying the extension of the same sharp rank-sensitive bound
beyond normal matrices as an open problem \cite{CostaRicoWolf}. Fu,
Gao, and Park recently resolved the rank-two endpoint through a sharp
estimate for the $S(2)$-norm of a double-antisymmetric projection
\cite{fu2026solution2copydistillabilitywerner}. Using that
projection estimate, the argument below obtains the arbitrary-matrix
inequality for every rank. Its main additional ingredient is a
completely positive decomposition of the negative partial-transpose
sector, followed by a complete-graph summation over the double-skew
subspace. The two-copy applications determine exact rank-constrained score
curves; the asymmetric formula yields a complete $r$-positivity
classification for the corresponding tensor-product maps. At local
dimension $d=r+1$, the critical signed filter becomes a four-qudit,
one-sided Schmidt-number witness with perfect completeness and exact
soundness $r/(r+1)$ against states of Schmidt number at most $r$.
For arbitrary copy number, we derive the inclusion--exclusion
polynomial, an exact code-state operator certificate, tensor
propagation, several exact regimes, and a universal block-positivity
radius. These results isolate, rather than solve, the collective
higher-copy Werner endpoint. We also expose the complete-graph Gram
defect that drives the rank-$r$ proof. Edgewise nonnegativity in that
decomposition is precisely the Fu--Gao--Park projection estimate.

The Werner-state and distillability terminology used in the applications
follows the standard conventions \cite{Werner1989,Horodecki1998}.
Throughout, $\norm{\cdot}_2$ denotes the Hilbert--Schmidt, or
Schatten-$2$, norm. The notation $\Tr_U$ means the partial trace
\emph{over} the factor $U$. Inner products are conjugate-linear in the
first argument and linear in the second. Relative to chosen orthonormal
bases, vectorization is normalized by
\[
    \operatorname{vec}(\ketbra{x}{y})
    =x\otimes\overline y.
\]

\begin{theorem}[Rank-$r$ partial-trace inequality]
\label[theorem]{thm:rank_r}
Let $U$ and $V$ be finite-dimensional complex Hilbert spaces. If
$C\in\cL(U\otimes V)$ is nonzero and $r_0:=\rank C$, then
\begin{equation}
\label{eq:main-bound-exact-rank}
    \norm{\Tr_U C}_2^2+\norm{\Tr_V C}_2^2
    \le r_0\norm{C}_2^2+\frac{1}{r_0}\abs{\Tr C}^2.
\end{equation}
Consequently, for every $r\in\mathbb N_{\ge1}$ such that $\rank C\le r$,
\begin{equation}
\label{eq:main-bound}
    \norm{\Tr_U C}_2^2+\norm{\Tr_V C}_2^2
    \le r\norm{C}_2^2+\frac{1}{r}\abs{\Tr C}^2.
\end{equation}
The assertion is immediate when $C=0$.
\end{theorem}

The coefficients are optimal whenever $\dim U\ge r$ and $\dim V\ge2$,
or with $U$ and $V$ interchanged. More precisely, the coefficient $r$
of $\norm{C}_2^2$ is optimal, and conditional on choosing this optimal
coefficient, the coefficient $1/r$ of $\abs{\Tr C}^2$ is optimal.

Indeed, choose
\[
    C=P_r\otimes\ketbra{v}{w},
\]
where $P_r$ is a rank-$r$ projection on $U$ and $v,w\in V$ are unit
vectors. Then
\[
\begin{aligned}
    \norm{\Tr_U C}_2^2&=r^2,\\
    \norm{\Tr_V C}_2^2&=r\abs{\ip{w}{v}}^2,\\
    \norm{C}_2^2&=r,\\
    \abs{\Tr C}^2&=r^2\abs{\ip{w}{v}}^2.
\end{aligned}
\]
Thus equality holds in \cref{eq:main-bound}. For a bound of the form
\[
    \norm{\Tr_U C}_2^2+\norm{\Tr_V C}_2^2
    \le a\norm{C}_2^2+b\abs{\Tr C}^2,
\]
taking $v\perp w$ forces $a\ge r$. Once $a=r$ is fixed, taking $v=w$
forces $b\ge1/r$.

\section{Double-Skew Estimates and Qudit Marginals}

We first record the double-skew Ky Fan estimate used below.

\begin{lemma}[Double-skew action bound]
\label[lemma]{lem:double-skew}
Fix orthonormal bases of $U$ and $V$, and let
\[
    \mathfrak{so}(U)
    :=\{L\in\cL(U):L^T=-L\},
    \qquad
    \mathfrak{so}(V)
    :=\{M\in\cL(V):M^T=-M\}.
\]
For every
\[
    K\in\mathfrak{so}(U)\otimes\mathfrak{so}(V),
\]
the two largest singular values satisfy
\[
    s_1(K)^2+s_2(K)^2\le\frac12\norm{K}_2^2.
\]
Consequently, for every orthonormal pair $x,y\in U\otimes V$,
\begin{equation}
\label{eq:double-skew-action}
    \norm{Kx}^2+\norm{Ky}^2
    \le\frac12\norm{K}_2^2.
\end{equation}
\end{lemma}

\begin{proof}
If $K=0$, both assertions are immediate. Hence assume $K\ne0$. Then
$\dim U,\dim V\ge2$. Let $d=\max\{\dim U,\dim V\}$. Embed $U$ and
$V$ isometrically into
$\C^d$ and extend every operator on a smaller factor by zero. The
resulting extension of $K$ belongs to
$\mathfrak{so}(\C^d)\otimes\mathfrak{so}(\C^d)$ and has the same
nonzero singular values and Hilbert--Schmidt norm as $K$. It is
therefore enough to prove the assertion when both local spaces have
dimension $d$.

Label the four vectorization factors as
\[
    U_{\mathrm{out}}\otimes V_{\mathrm{out}}
    \otimes U_{\mathrm{in}}\otimes V_{\mathrm{in}},
\]
where the chosen bases identify each input conjugate space with its
corresponding output space. The Schmidt bipartition used below is
\[
    (U_{\mathrm{out}}V_{\mathrm{out}})
    :
    (U_{\mathrm{in}}V_{\mathrm{in}}).
\]
Since $L^T=-L$ implies that $\operatorname{vec}(L)$ is antisymmetric
between its output and input copies, vectorization of
$K\in\mathfrak{so}(U)\otimes\mathfrak{so}(V)$ gives
\[
    Q_-\operatorname{vec}(K)=\operatorname{vec}(K),
\]
where, on the displayed four factors,
\[
    Q_-
    :=P_{\wedge^2 U}^{(U_{\mathrm{out}},U_{\mathrm{in}})}
      \otimes
      P_{\wedge^2 V}^{(V_{\mathrm{out}},V_{\mathrm{in}})}.
\]
Equivalently, after the permutation that groups the two copies of $U$
and the two copies of $V$,
$\operatorname{vec}(K)\in(\wedge^2U)\otimes(\wedge^2V)$, while the
Schmidt cut remains the crossed output--input cut above.

Set
\[
    \phi:=\frac{\operatorname{vec}(K)}{\norm{K}_2}.
\]
The Schmidt coefficients of $\phi$ across the output--input
bipartition are precisely $s_i(K)/\norm{K}_2$. Hence the best
Schmidt-rank-two approximation identity gives
\begin{equation}
\label{eq:best-SR2-approximation}
    \max_{\substack{\norm{z}=1\\ \SR(z)\le2}}
    \abs{\ip{z}{\phi}}^2
    =
    \frac{s_1(K)^2+s_2(K)^2}{\norm{K}_2^2};
\end{equation}
see, for example, \cite[Theorem 3.3]{JohnstonKribs2010}.
For every unit vector $z$ of Schmidt rank at most $2$,
\[
    \abs{\ip{z}{\phi}}^2
    =\abs{\ip{Q_-z}{\phi}}^2
    \le \norm{Q_-z}^2
    =\ip{z}{Q_-z}.
\]
The sharp double-antisymmetric projection estimate of Fu, Gao, and
Park \cite[Theorem 2.4]{fu2026solution2copydistillabilitywerner},
with antisymmetric pairs
$(U_{\mathrm{out}},U_{\mathrm{in}})$ and
$(V_{\mathrm{out}},V_{\mathrm{in}})$ and Schmidt cut
$(U_{\mathrm{out}}V_{\mathrm{out}}):(U_{\mathrm{in}}V_{\mathrm{in}})$,
now gives
\[
    \ip{z}{Q_-z}\le\frac12.
\]
Combining this with \cref{eq:best-SR2-approximation} yields
\[
    s_1(K)^2+s_2(K)^2\le\frac12\norm{K}_2^2.
\]

Finally, the Ky Fan maximum principle gives
\[
    \norm{Kx}^2+\norm{Ky}^2
    \le s_1(K)^2+s_2(K)^2
\]
for every orthonormal pair $x,y$, proving
\cref{eq:double-skew-action}.
\end{proof}

We next establish an operator inequality for pure states with a
maximally mixed marginal.

\begin{proposition}[Rank-$r$ marginal operator inequality]
\label[proposition]{prop:operator_ineq}
Let $Q=\C^r$, and let $\psi\in U\otimes V\otimes Q$ be a unit vector
such that
\[
    \rho_Q
    :=\Tr_{UV}\ketbra{\psi}{\psi}
    =\frac1r I_Q.
\]
Writing
\[
    \rho=\ketbra{\psi}{\psi},
    \qquad
    \rho_{UQ}=\Tr_V\rho,
    \qquad
    \rho_{VQ}=\Tr_U\rho,
\]
one has
\begin{equation}
\label{eq:H_r}
    H_r
    :=
    I_{UVQ}+\frac1r\rho
    -\rho_{UQ}\otimes I_V
    -I_U\otimes\rho_{VQ}
    \succeq0.
\end{equation}
Here and below, operators are placed on their labeled tensor factors,
regardless of whether those factors are adjacent in the displayed
tensor ordering.
\end{proposition}

\begin{proof}
Fix orthonormal bases of $U$ and $V$, relative to which transposition
and the skew-symmetric operator spaces are defined. Although the proof
uses these bases, the resulting operator inequality is basis
independent. By the maximally mixed condition on $Q$, the vector $\psi$
has a Schmidt decomposition
\begin{equation}
\label{eq:psi-schmidt}
    \psi
    =
    \frac1{\sqrt r}
    \sum_{i=1}^r e_i\otimes q_i,
\end{equation}
where $\{e_i\}_{i=1}^r$ is an orthonormal family in $U\otimes V$ and
$\{q_i\}_{i=1}^r$ is an orthonormal basis of $Q$.

Expanding $\rho$ and transposing the $UV$ factor gives
\begin{align}
    \rho^{T_{UV}}
    &=
    \frac1r
    \sum_{i,j=1}^r
    \bigl(\ketbra{e_i}{e_j}\bigr)^T
    \otimes\ketbra{q_i}{q_j}
    \notag\\
    &=
    \frac1r
    \sum_{i,j=1}^r
    \ketbra{\overline e_j\otimes q_i}
            {\overline e_i\otimes q_j}.
\label{eq:rho-partial-transpose-expanded}
\end{align}
Let
\[
    S:=\operatorname{span}
    \{\overline e_1,\ldots,\overline e_r\},
\]
and let $J:Q\to S$ be the unitary determined by
$Jq_i=\overline e_i$. Define a self-adjoint unitary $F_J$ on
$S\otimes Q$ by
\[
    F_J(s\otimes q)
    :=Jq\otimes J^{-1}s.
\]
In particular,
\[
    F_J(\overline e_i\otimes q_j)
    =\overline e_j\otimes q_i.
\]
Thus the sum in \cref{eq:rho-partial-transpose-expanded} is $F_J$ on
$S\otimes Q$. Let
\[
    P_\pm:=\frac12(I_{S\otimes Q}\pm F_J)
\]
be its positive and negative eigenspace projections, extended by zero
on $(S\otimes Q)^\perp$. Then
\begin{equation}
\label{eq:rho-partial-transpose}
    \rho^{T_{UV}}
    =
    \frac1r(P_+-P_-).
\end{equation}
Explicitly,
\[
    P_-
    =
    \sum_{1\le i<j\le r}
    \ketbra{\eta_{ij}}{\eta_{ij}},
\]
where
\begin{equation}
\label{eq:eta-ij}
    \eta_{ij}
    =
    \frac1{\sqrt2}
    \bigl(
      \overline e_i\otimes q_j
      -
      \overline e_j\otimes q_i
    \bigr).
\end{equation}

For a finite-dimensional Hilbert space $E$, define
\[
    \Lambda_E(X)
    :=
    \Tr(X)I_E-X^T.
\]
If $L_{ab}=E_{ab}-E_{ba}$, then
\begin{equation}
\label{eq:Lambda-kraus}
    \Lambda_E(X)
    =
    \sum_{a<b}L_{ab}XL_{ab}^*.
\end{equation}
Thus $\Lambda_E$ is completely positive. Equivalently, its Choi
operator is
\[
    J(\Lambda_E)
    =
    \sum_{i,j}
    \ketbra{i}{j}\otimes\Lambda_E(E_{ij})
    =
    I-F,
\]
whose spectrum is contained in $\{0,2\}$.

Define
\[
    \Phi
    :=
    \Lambda_U\otimes\Lambda_V\otimes\id_Q.
\]
For an operator $Y$ on $U\otimes V\otimes Q$, direct expansion gives
\begin{equation}
\label{eq:Phi-expansion}
\begin{aligned}
    \Phi(Y)
    ={}&
    I_{UV}\otimes\Tr_{UV}Y
    -I_U\otimes(\Tr_UY)^{T_V}\\
    &-(\Tr_VY)^{T_U}\otimes I_V
    +Y^{T_{UV}},
\end{aligned}
\end{equation}
with all terms placed on their labeled tensor factors.

Applying \cref{eq:Phi-expansion} to
$Y=\rho^{T_{UV}}$ gives
\begin{equation}
\label{eq:Phi-rhoT}
    \Phi(\rho^{T_{UV}})
    =
    \frac1r I_{UVQ}
    -\rho_{UQ}\otimes I_V
    -I_U\otimes\rho_{VQ}
    +\rho.
\end{equation}
Using \cref{eq:rho-partial-transpose}, we therefore obtain
\begin{equation}
\label{eq:H-decomposition}
    H_r
    =
    \frac1r\Phi(P_+)
    +
    \frac1r
    \Bigl[
      (r-1)(I_{UVQ}-\rho)-\Phi(P_-)
    \Bigr].
\end{equation}
Since $\Phi(P_+)\succeq0$, it remains to prove
\begin{equation}
\label{eq:Phi-Pminus-goal}
    \Phi(P_-)
    \preceq
    (r-1)(I_{UVQ}-\rho).
\end{equation}

Let $\{L_\alpha\}$ and $\{M_\beta\}$ be the standard skew-symmetric
Kraus operators for $\Lambda_U$ and $\Lambda_V$, normalized by
\[
    \langle L_\alpha,L_{\alpha'}\rangle_{\mathrm{HS}}
    =
    2\delta_{\alpha\alpha'},
    \qquad
    \langle M_\beta,M_{\beta'}\rangle_{\mathrm{HS}}
    =
    2\delta_{\beta\beta'}.
\]
Define a linear map $\mathcal T$ on coefficient arrays
$c=(c_{\alpha\beta}^{(ij)})$ by
\begin{equation}
\label{eq:T-definition}
    \mathcal Tc
    :=
    \sum_{1\le i<j\le r}
    \sum_{\alpha,\beta}
    c_{\alpha\beta}^{(ij)}
    (L_\alpha\otimes M_\beta\otimes I_Q)\eta_{ij}.
\end{equation}
Indeed, using the Kraus form of $\Phi$ and the rank-one expansion of
$P_-$,
\begin{align}
    \Phi(P_-)
    &=
    \sum_{i<j}\sum_{\alpha,\beta}
    \ketbra{
      (L_\alpha\otimes M_\beta\otimes I_Q)\eta_{ij}}
      {(L_\alpha\otimes M_\beta\otimes I_Q)\eta_{ij}}
    \notag\\
    &=\mathcal T\mathcal T^*.
\label{eq:Phi-Pminus-TTstar}
\end{align}
The last equality follows because the vectors in the first line are
exactly the columns of $\mathcal T$ indexed by
$(i,j,\alpha,\beta)$.

For each edge $(i,j)$ of the complete graph $K_r$, define
\[
    K_c^{(ij)}
    :=
    \sum_{\alpha,\beta}
    c_{\alpha\beta}^{(ij)}
    L_\alpha\otimes M_\beta.
\]
Then
\[
    K_c^{(ij)}
    \in
    \mathfrak{so}(U)\otimes\mathfrak{so}(V),
\]
and orthogonality of the Kraus operators gives
\begin{equation}
\label{eq:Kc-norm}
\begin{aligned}
    \norm{K_c^{(ij)}}_2^2
    &=
    \sum_{\alpha,\beta}
    \abs{c_{\alpha\beta}^{(ij)}}^2
    \norm{L_\alpha\otimes M_\beta}_2^2\\
    &=
    4\norm{c^{(ij)}}_2^2.
\end{aligned}
\end{equation}

Grouping the output of $\mathcal T$ according to the basis
$\{q_k\}$ gives
\begin{equation}
\label{eq:T-by-vertices}
    \mathcal Tc
    =
    \frac1{\sqrt2}
    \sum_{k=1}^r
    \left(
      \sum_{i<k}K_c^{(ik)}\overline e_i
      -
      \sum_{k<j}K_c^{(kj)}\overline e_j
    \right)
    \otimes q_k.
\end{equation}
At each vertex $k$, there are $r-1$ incident terms. Hence
Cauchy--Schwarz yields
\begin{align}
    \norm{\mathcal Tc}^2
    &\le
    \frac{r-1}{2}
    \sum_{k=1}^r
    \left(
      \sum_{i<k}
      \norm{K_c^{(ik)}\overline e_i}^2
      +
      \sum_{k<j}
      \norm{K_c^{(kj)}\overline e_j}^2
    \right)
    \notag\\
    &=
    \frac{r-1}{2}
    \sum_{1\le i<j\le r}
    \left(
      \norm{K_c^{(ij)}\overline e_i}^2
      +
      \norm{K_c^{(ij)}\overline e_j}^2
    \right).
\label{eq:T-Cauchy}
\end{align}
By \cref{lem:double-skew}, each parenthesized expression is at most
$\frac12\norm{K_c^{(ij)}}_2^2$. Therefore,
\begin{align}
    \norm{\mathcal Tc}^2
    &\le
    \frac{r-1}{4}
    \sum_{i<j}
    \norm{K_c^{(ij)}}_2^2
    \notag\\
    &=
    (r-1)\norm{c}_2^2,
\label{eq:T-operator-bound}
\end{align}
where the last equality follows from \cref{eq:Kc-norm}. Thus
\[
    \mathcal T\mathcal T^*
    \preceq
    (r-1)I_{UVQ}.
\]

It remains to use the orthogonality of the range of $\mathcal T$ to
$\psi$. Since each tensor product of two skew-symmetric operators is
symmetric,
\[
    K^T=K
    \qquad
    \text{for every }
    K\in
    \mathfrak{so}(U)\otimes\mathfrak{so}(V).
\]
Consequently,
\begin{align}
    \ip{\psi}{\mathcal Tc}
    &=
    \frac1{\sqrt{2r}}
    \sum_{i<j}
    \left(
      \overline e_j^T
      K_c^{(ij)}
      \overline e_i
      -
      \overline e_i^T
      K_c^{(ij)}
      \overline e_j
    \right)
    \notag\\
    &=0.
\label{eq:T-orthogonal}
\end{align}
Hence
\[
    \operatorname{ran}\mathcal T
    \subseteq
    \psi^\perp.
\]
Since $\ran\mathcal T\subseteq\psi^\perp$, one has
\[
    \mathcal T\mathcal T^*
    =P_{\psi^\perp}\mathcal T\mathcal T^*P_{\psi^\perp}.
\]
Moreover, \cref{eq:T-operator-bound} says that
$\norm{\mathcal T}^2\le r-1$. Therefore, for every vector $x$,
\[
\begin{aligned}
    \ip{x}{\mathcal T\mathcal T^*x}
    &=
    \ip{P_{\psi^\perp}x}
       {\mathcal T\mathcal T^*P_{\psi^\perp}x}\\
    &\le
    (r-1)\norm{P_{\psi^\perp}x}^2.
\end{aligned}
\]
Equivalently,
\[
    \mathcal T\mathcal T^*
    \preceq
    (r-1)P_{\psi^\perp}
    =
    (r-1)(I_{UVQ}-\rho).
\]
This proves \cref{eq:Phi-Pminus-goal}, and therefore
$H_r\succeq0$ by \cref{eq:H-decomposition}.
\end{proof}

\section{Proof of the Rank-\texorpdfstring{$r$}{r}
Partial-Trace Inequality}

\begin{lemma}[Partial-trace contraction identities]
\label[lemma]{lem:partial-trace-contractions}
Let $e_1,\ldots,e_s\in U\otimes V$ be orthonormal, let
$d_1,\ldots,d_s\in U\otimes V$, and let $q_1,\ldots,q_s$ be an
orthonormal basis of $Q=\C^s$. Define
\[
    C:=\sum_{i=1}^s\ketbra{e_i}{d_i},
    \qquad
    \psi:=\frac1{\sqrt s}\sum_{i=1}^s e_i\otimes q_i,
    \qquad
    \rho:=\ketbra{\psi}{\psi},
\]
and
\[
    \delta:=\sum_{i=1}^s d_i\otimes q_i.
\]
Writing $\rho_{UQ}=\Tr_V\rho$ and $\rho_{VQ}=\Tr_U\rho$, one has
\begin{align}
    \ip{\delta}{\delta}
    &=\norm{C}_2^2,
\label{eq:contraction-norm}\\
    \ip{\delta}{\rho\delta}
    &=\frac1s\abs{\Tr C}^2,
\label{eq:contraction-trace}\\
    \ip{\delta}{(\rho_{UQ}\otimes I_V)\delta}
    &=\frac1s\norm{\Tr_U C}_2^2,
\label{eq:contraction-U}\\
    \ip{\delta}{(I_U\otimes\rho_{VQ})\delta}
    &=\frac1s\norm{\Tr_V C}_2^2.
\label{eq:contraction-V}
\end{align}
\end{lemma}

\begin{proof}
The first identity follows from the orthonormality of the $e_i$:
\[
    \norm{C}_2^2
    =\sum_i\norm{d_i}^2
    =\ip{\delta}{\delta}.
\]
Since $\Tr C=\sum_i\ip{d_i}{e_i}$ under the stated inner-product
convention,
\[
    \ip{\psi}{\delta}
    =\frac1{\sqrt s}\sum_i\ip{e_i}{d_i}
    =\frac1{\sqrt s}\overline{\Tr C},
\]
which proves \cref{eq:contraction-trace}.

For the first marginal contraction,
\[
    \rho_{UQ}
    =\frac1s\sum_{i,j}
      \Tr_V\ketbra{e_i}{e_j}\otimes\ketbra{q_i}{q_j},
\]
and hence
\begin{equation}
\label{eq:contraction-U-expansion}
    \ip{\delta}{(\rho_{UQ}\otimes I_V)\delta}
    =\frac1s\sum_{i,j}
      \ip{d_i}{
        (\Tr_V\ketbra{e_i}{e_j}\otimes I_V)d_j
      }.
\end{equation}
Choose product bases $\{u_a\}$ of $U$ and $\{v_b\}$ of $V$, and
write
\[
    e_i=\sum_{a,b}e_{i,ab}u_a\otimes v_b,
    \qquad
    d_i=\sum_{a,b}d_{i,ab}u_a\otimes v_b.
\]
Then
\[
    (\Tr_UC)_{b,c}
    =\sum_{i,a}e_{i,ab}\overline{d_{i,ac}},
\]
so
\[
    \norm{\Tr_UC}_2^2
    =\sum_{i,j}\sum_{a,a'}\sum_{b,c}
      e_{i,ab}\overline{d_{i,ac}}
      \overline{e_{j,a'b}}d_{j,a'c}.
\]
Expanding the right-hand side of
\cref{eq:contraction-U-expansion} gives the same sum, proving
\cref{eq:contraction-U}. Interchanging $U$ and $V$ proves
\cref{eq:contraction-V}.
\end{proof}

\begin{proof}[Proof of \cref{thm:rank_r}]
The case $C=0$ is immediate. Suppose $C\ne0$, and let
\[
    r_0:=\rank C.
\]
Choose an orthonormal basis
$\{e_i\}_{i=1}^{r_0}$ of $\operatorname{ran}C$ and set
\[
    d_i:=C^*e_i.
\]
Then
\begin{equation}
\label{eq:C-decomposition}
    C
    =
    \sum_{i=1}^{r_0}\ketbra{e_i}{d_i},
    \qquad
    \sum_{i=1}^{r_0}\norm{d_i}^2
    =
    \norm{C}_2^2.
\end{equation}
Since $\rank C\le r$, one has $1\le r_0\le r$.

Let $Q=\C^{r_0}$, with orthonormal basis
$\{q_i\}_{i=1}^{r_0}$, and define
\[
    \psi
    =
    \frac1{\sqrt{r_0}}
    \sum_{i=1}^{r_0}e_i\otimes q_i,
    \qquad
    \rho=\ketbra{\psi}{\psi}.
\]
Then
\[
    \rho_Q
    =
    \Tr_{UV}\rho
    =
    \frac1{r_0}I_Q.
\]
Also define the unnormalized vector
\[
    \delta
    :=
    \sum_{i=1}^{r_0}d_i\otimes q_i
    \in U\otimes V\otimes Q.
\]

Applying \cref{prop:operator_ineq} with $r=r_0$ gives
$H_{r_0}\succeq0$. Therefore,
\begin{align}
    0
    &\le
    \ip{\delta}{H_{r_0}\delta}
    \notag\\
    &=
    \ip{\delta}{\delta}
    +
    \frac1{r_0}\ip{\delta}{\rho\delta}
    -
    \ip{\delta}{
      (\rho_{UQ}\otimes I_V)\delta
    }
    -
    \ip{\delta}{
      (I_U\otimes\rho_{VQ})\delta
    }.
\label{eq:H-expectation}
\end{align}
By \cref{lem:partial-trace-contractions},
\begin{align*}
    \ip{\delta}{\delta}
    &=\norm{C}_2^2,\\
    \ip{\delta}{\rho\delta}
    &=\frac1{r_0}\abs{\Tr C}^2,\\
    \ip{\delta}{(\rho_{UQ}\otimes I_V)\delta}
    &=\frac1{r_0}\norm{\Tr_U C}_2^2,\\
    \ip{\delta}{(I_U\otimes\rho_{VQ})\delta}
    &=\frac1{r_0}\norm{\Tr_V C}_2^2.
\end{align*}
Substituting these identities into \cref{eq:H-expectation} and
multiplying by $r_0$ gives
\begin{equation}
\label{eq:main-bound-r0}
    \norm{\Tr_U C}_2^2+\norm{\Tr_V C}_2^2
    \le
    r_0\norm{C}_2^2
    +
    \frac1{r_0}\abs{\Tr C}^2.
\end{equation}

It remains to pass from $r_0$ to $r$. Let $P$ be the orthogonal
projection onto $\ran C$. Since $PC=C$ and $\norm{P}_2^2=r_0$,
Hilbert--Schmidt Cauchy--Schwarz gives the trace-rank estimate
\[
    \abs{\Tr C}^2
    =\abs{\Tr(PC)}^2
    \le
    \norm{P}_2^2\norm{C}_2^2
    =
    r_0\norm{C}_2^2.
\]
Thus
\begin{align}
    &
    r\norm{C}_2^2
    +
    \frac1r\abs{\Tr C}^2
    -
    \left(
      r_0\norm{C}_2^2
      +
      \frac1{r_0}\abs{\Tr C}^2
    \right)
    \notag\\
    &\qquad
    =
    (r-r_0)\norm{C}_2^2
    -
    \frac{r-r_0}{rr_0}\abs{\Tr C}^2
    \notag\\
    &\qquad
    \ge
    (r-r_0)\norm{C}_2^2
    -
    \frac{r-r_0}{r}\norm{C}_2^2
    \notag\\
    &\qquad
    \ge0.
\label{eq:rank-monotonicity}
\end{align}
Combining \cref{eq:main-bound-r0} and
\cref{eq:rank-monotonicity} proves \cref{eq:main-bound}.
\end{proof}

\section{Gram Structure and Defect Identities}

This section records the algebraic defect structure of the rank-$r$
proof. The main residual is represented by positive Gram operators
once the range isometry is fixed; this is not claimed to be a uniform
bounded-degree polynomial sum-of-squares identity in the raw Stiefel
coordinates. The only literal polynomial SOS below is the elementary
trace--rank residual.

For an operator $C$ of rank at most $r$, define the two nonnegative
rank-constrained residuals
\begin{align}
    \mathsf G_r(C)
    &:=%
    r\norm C_2^2+\frac1r\abs{\Tr C}^2
    -\norm{\Tr_U C}_2^2-\norm{\Tr_V C}_2^2,
\label{eq:sos-main-gap}\\
    \mathsf R_r(C)
    &:=%
    r\norm C_2^2-\abs{\Tr C}^2.
\label{eq:sos-trace-rank-gap}
\end{align}
The first is the main theorem's deficit and the second is the ordinary
trace--rank Cauchy--Schwarz deficit.

\begin{proposition}[Polynomial SOS for the trace--rank residual]
\label[proposition]{prop:sos-trace-rank}
Assume $r\le\dim(U\otimes V)$ and write
\[
    C=\sum_{i=1}^r\ketbra{e_i}{d_i},
\]
where $e_1,\ldots,e_r$ are orthonormal and zero columns $d_i=0$ are
allowed.  In the Hilbert space
$(U\otimes V)\otimes\C^r$, put
\[
    \mathbf e:=\sum_i e_i\otimes q_i,
    \qquad
    \mathbf d:=\sum_i d_i\otimes q_i.
\]
Equip the exterior square with the convention that
$u_\alpha\wedge u_\beta$ for $\alpha<\beta$ is orthonormal
whenever $(u_\alpha)$ is an orthonormal basis. Then
\begin{equation}
\label{eq:sos-lagrange}
    \mathsf R_r(C)
    =\norm{\mathbf e\wedge\mathbf d}^2.
\end{equation}
Equivalently, in any orthonormal coordinate system,
\begin{equation}
\label{eq:sos-lagrange-coordinates}
    \mathsf R_r(C)
    =
    \sum_{\alpha<\beta}
    \abs{
      \mathbf e_\alpha\mathbf d_\beta
      -\mathbf e_\beta\mathbf d_\alpha
    }^2.
\end{equation}
Hence $\mathsf R_r$ has a degree-two complex, or degree-four real,
polynomial SOS certificate on the range-factorization variables.
\end{proposition}

\begin{proof}
One has
\[
    \norm{\mathbf e}^2=r,
    \qquad
    \norm{\mathbf d}^2=\norm C_2^2,
    \qquad
    \abs{\ip{\mathbf e}{\mathbf d}}
    =\abs{\Tr C}.
\]
The complex Lagrange identity
\[
    \norm{a}^2\norm{b}^2-\abs{\ip{a}{b}}^2
    =\sum_{\alpha<\beta}
      \abs{a_\alpha b_\beta-a_\beta b_\alpha}^2
\]
therefore gives both displayed formulas.
\end{proof}

For the remainder of this section assume
$r\le\dim(U\otimes V)$, so that an $r$-column range isometry is
available; operators of smaller rank are represented by allowing zero
columns.

We next extract the Gram certificate implicit in the proof of
\cref{thm:rank_r}.  Retain an exact rank-$r$ range factorization
\[
    C=\sum_{i=1}^r\ketbra{e_i}{d_i},
    \qquad
    \delta=\sum_i d_i\otimes q_i,
\]
and define the range isometry
\[
    E:\C^r\longrightarrow U\otimes V,
    \qquad Eq_i=e_i.
\]
Retain also the associated state $\psi$, projection $\rho$, twisted
symmetric and antisymmetric projections $P_\pm$, and completely
positive map $\Phi$ from the proof of
\cref{prop:operator_ineq}.  Let $\mathcal T_-$ be the synthesis map
in \cref{eq:T-definition}.  Choose an orthonormal basis of
$\ran P_+$ and let $\mathcal T_+$ be the synthesis map whose columns
are all Kraus images under $\Phi$ of those positive-sector basis
vectors.  Then
\[
    \Phi(P_\pm)=\mathcal T_\pm\mathcal T_\pm^*.
\]

\begin{proposition}[Exact positive/negative-sector Gram identity]
\label[proposition]{prop:sos-sector-gram}
Put
\begin{equation}
\label{eq:sos-negative-defect}
    B_E
    :=
    (r-1)(I-\rho)-\mathcal T_-\mathcal T_-^*.
\end{equation}
Then
\begin{equation}
\label{eq:sos-main-gram}
    \boxed{
    \mathsf G_r(C)
    =
    \norm{\mathcal T_+^*\delta}^2
    +\ip{\delta}{B_E\delta}.}
\end{equation}
Moreover, $B_E\succeq0$.
\end{proposition}

\begin{proof}
Multiplying \cref{eq:H-decomposition} by $r$ and using
\cref{lem:partial-trace-contractions} gives
\[
\begin{aligned}
    \mathsf G_r(C)
    &=%
    \ip{\delta}{\Phi(P_+)\delta}\\
    &\quad+
    \ip{\delta}{
      \left[(r-1)(I-\rho)-\Phi(P_-)\right]
      \delta}.
\end{aligned}
\]
Substitution of the two synthesis factorizations proves
\cref{eq:sos-main-gram}.  The positivity of $B_E$ follows from the
complete-graph defect certificate below together with
$\ran\mathcal T_-\subseteq\psi^\perp$.
\end{proof}

For a coefficient array $c=(c^{(ij)}_{\alpha\beta})$, define
\[
    K_{ij}
    :=
    \sum_{\alpha,\beta}
    c^{(ij)}_{\alpha\beta}L_\alpha\otimes M_\beta.
\]
At a vertex $k$, associate to an incident edge the signed vectors
\[
    v_{k,ik}:=K_{ik}\overline e_i\quad(i<k),
    \qquad
    v_{k,kj}:=-K_{kj}\overline e_j\quad(k<j).
\]

\begin{proposition}[Complete-graph defect certificate]
\label[proposition]{prop:sos-complete-graph}
With the skew Kraus normalization used in
\cref{eq:Kc-norm},
\begin{equation}
\label{eq:sos-complete-graph}
\boxed{
\begin{aligned}
 &(r-1)\norm c^2-\norm{\mathcal T_-c}^2\\
 &\quad=
 \frac{r-1}{2}\sum_{i<j}
 \left[
   \frac12\norm{K_{ij}}_2^2
   -\norm{K_{ij}\overline e_i}^2
   -\norm{K_{ij}\overline e_j}^2
 \right]\\
 &\qquad+
 \frac12\sum_{k=1}^r
 \sum_{\substack{e<f\\e,f\ni k}}
 \norm{v_{k,e}-v_{k,f}}^2.
\end{aligned}}
\end{equation}
The vertex-variance terms are manifestly nonnegative, and the edge
brackets are nonnegative by the Fu--Gao--Park estimate in
\cref{lem:double-skew}. Consequently
\[
    \norm{\mathcal T_-}^2\le r-1,
\]
and, since its range is orthogonal to $\psi$,
$B_E\succeq0$.
\end{proposition}

\begin{proof}
At every vertex use the exact variance identity
\[
    (r-1)\sum_{e\ni k}\norm{v_{k,e}}^2
    -\norm{\sum_{e\ni k}v_{k,e}}^2
    =
    \sum_{\substack{e<f\\e,f\ni k}}
    \norm{v_{k,e}-v_{k,f}}^2.
\]
Sum over vertices and use
\cref{eq:T-by-vertices}.  Each edge contributes its two endpoint
actions, while
$\norm{K_{ij}}_2^2=4\norm{c^{(ij)}}^2$ by
\cref{eq:Kc-norm}.  Adding and subtracting the total edge action gives
exactly \cref{eq:sos-complete-graph}.  The edge brackets are
nonnegative by \cref{lem:double-skew}.  The final operator conclusion
then follows exactly as in the last paragraph of the proof of
\cref{prop:operator_ineq}.
\end{proof}

\begin{remark}[Scope of the Gram certificate]
The complete-graph decomposition \cref{eq:sos-complete-graph}
separates the global synthesis-map defect into vertex variances and
edgewise double-skew deficits. Edgewise nonnegativity is precisely the
Fu--Gao--Park projection estimate used in \cref{lem:double-skew}.
No uniform polynomial SOS certificate in the raw range-isometry
variables is claimed.
\end{remark}

\section{Two-Copy Filters and Positive Maps}

This section develops the two-copy consequences closest to the
main inequality: the exact symmetric score curve, an asymmetric
score theorem, and the resulting block-positivity and positive-map
thresholds.

\subsection{Two-copy block positivity}

Fix $d\ge2$ and $\alpha\in\mathbb R$. Let
\[
    \rho_\alpha
    :=
    I+\alpha F
\]
be the unnormalized Werner operator on $\C^d\otimes\C^d$; for
$\alpha\in[-1,1]$, it is proportional to the standard Werner state
\cite{Werner1989}. Its partial transpose is
\[
    \rho_\alpha^\Gamma
    =
    I+\alpha\ketbra{\Omega}{\Omega},
    \qquad
    \Omega
    :=
    \sum_{i=1}^d e_i\otimes e_i.
\]
The case $\alpha\ge0$ is immediate because
\[
    I+\alpha\ketbra{\Omega}{\Omega}\succeq0.
\]

For $t\ge0$, define
\begin{equation}
\label{eq:q2}
    q^{(2)}(-t,C)
    :=
    \norm{C}_2^2
    -
    t\left(
      \norm{\Tr_1C}_2^2+\norm{\Tr_2C}_2^2
    \right)
    +
    t^2\abs{\Tr C}^2.
\end{equation}

An operator is called $r$-block-positive if its expectation is
nonnegative on every vector of Schmidt rank at most $r$. With the
vectorization convention fixed in the introduction and the regrouping
$(A_1A_2):(B_1B_2)$, one has
\begin{equation}
\label{eq:q2-vectorization}
\begin{aligned}
    &\left\langle
      \operatorname{vec}(C),
      \bigl(I-t\ketbra{\Omega}{\Omega}\bigr)^{\otimes2}
      \operatorname{vec}(C)
    \right\rangle\\
    &\qquad =q^{(2)}(-t,C).
\end{aligned}
\end{equation}
Moreover,
\[
    \SR\bigl(\operatorname{vec}(C)\bigr)=\rank C
\]
across the same output--input regrouping. Thus $r$-block positivity of
the two-copy operator is equivalent to $q^{(2)}(-t,C)\ge0$ for every
matrix $C$ of rank at most $r$.
For every $C$ of rank at most $r$, \cref{eq:main-bound} implies
\begin{align}
    q^{(2)}(-t,C)
    &\ge
    (1-rt)\norm{C}_2^2
    +
    \left(
      t^2-\frac tr
    \right)\abs{\Tr C}^2.
\label{eq:q2-first-bound}
\end{align}
If $0\le t\le1/r$, then
$t^2-t/r\le0$. Using
\[
    \abs{\Tr C}^2
    \le
    r\norm{C}_2^2
\]
in the appropriate direction gives
\begin{align}
    q^{(2)}(-t,C)
    &\ge
    \left(
      1-rt+rt^2-t
    \right)\norm{C}_2^2
    \notag\\
    &=
    (1-t)(1-rt)\norm{C}_2^2
    \ge0.
\label{eq:q2-final-bound}
\end{align}

\begin{proposition}[Exact rank-constrained two-copy score curve]
\label[proposition]{prop:exact-q2-score}
Let $d\ge2$ and $1\le r\le d$.  For $t\ge0$, define
\[
    \mu_r(t)
    :=
    \inf_{\substack{C\in M_{d^2}(\C)\setminus\{0\}\\
                     \rank C\le r}}
    \frac{q^{(2)}(-t,C)}{\norm{C}_2^2}.
\]
Then
\begin{equation}
\label{eq:exact-q2-score}
    \mu_r(t)
    =
    \begin{cases}
      (1-t)(1-rt),&0\le t\le 1/r,\\[1mm]
      1-rt,&t\ge1/r.
    \end{cases}
\end{equation}
Both branches are attained by rank-$r$ matrices.
\end{proposition}

\begin{proof}
For $C\ne0$ of rank at most $r$, set
\[
    z:=\frac{\abs{\Tr C}^2}{\norm{C}_2^2}.
\]
The trace--rank estimate gives $0\le z\le r$.  Dividing
\cref{eq:q2-first-bound} by $\norm{C}_2^2$ yields
\begin{equation}
\label{eq:q2-score-z}
    \frac{q^{(2)}(-t,C)}{\norm{C}_2^2}
    \ge
    1-rt+\left(t^2-\frac tr\right)z.
\end{equation}
If $0\le t\le1/r$, the coefficient of $z$ is nonpositive, so
\cref{eq:q2-score-z} and $z\le r$ give
\[
    \frac{q^{(2)}(-t,C)}{\norm{C}_2^2}
    \ge
    1-rt+r\left(t^2-\frac tr\right)
    =(1-t)(1-rt).
\]
Choose a rank-$r$ projection $P_r\in M_d(\C)$ and a unit vector
$v\in\C^d$.  For
\[
    C=P_r\otimes\ketbra{v}{v}
\]
one has
\[
\begin{aligned}
    \norm{C}_2^2&=r,
    &\abs{\Tr C}^2&=r^2,\\
    \norm{\Tr_1C}_2^2&=r^2,
    &\norm{\Tr_2C}_2^2&=r,
\end{aligned}
\]
and hence equality holds.

If $t\ge1/r$, the coefficient of $z$ in
\cref{eq:q2-score-z} is nonnegative, so
\[
    \frac{q^{(2)}(-t,C)}{\norm{C}_2^2}
    \ge1-rt.
\]
Choose orthogonal unit vectors $v,w\in\C^d$ and set
\[
    C=P_r\otimes\ketbra{v}{w}.
\]
Then
\[
    \Tr C=0,
    \qquad
    \norm{\Tr_1C}_2^2=r^2,
    \qquad
    \Tr_2C=0,
\]
so equality again holds.
\end{proof}

\begin{proposition}[Exact affine decomposition of the two-copy score]
\label[proposition]{prop:sos-affine-score}
For every $0\le t\le1/r$ and every $C$ of rank at most $r$,
\begin{equation}
\label{eq:sos-affine-score}
\begin{aligned}
    &q^{(2)}(-t,C)
    -(1-t)(1-rt)\norm C_2^2\\
    &\qquad=
    t\,\mathsf G_r(C)
    +t\left(\frac1r-t\right)\mathsf R_r(C).
\end{aligned}
\end{equation}
Thus nonnegativity of the two residuals certifies the two-copy lower
bound. It certifies the exact soundness curve when the dimension
hypotheses of \cref{prop:exact-q2-score} ensure that the baseline is
attained.
\end{proposition}

\begin{proof}
Substitute the definition of $\mathsf G_r(C)$ into
\cref{eq:q2} and collect terms.  The coefficient of $\norm C_2^2$ on
the right is
\[
    t r+t\left(\frac1r-t\right)r
    =t(r+1)-rt^2,
\]
and the coefficient of $\abs{\Tr C}^2$ is
\[
    \frac{t}{r}-t\left(\frac1r-t\right)=t^2.
\]
After adding the baseline
$(1-t)(1-rt)\norm C_2^2$, this is exactly
\cref{eq:q2}.
\end{proof}

\begin{remark}[Why the range $r\le d$ is essential]
The extremizers in \cref{prop:exact-q2-score} require a rank-$r$
projection in one $d$-dimensional tensor factor.  The score formula does
not in general extend to $r>d$ by replacing $r$ with
$\min\{r,d\}$.  For example, when the rank bound is $r=d^2$, no rank
constraint remains.  If $0\le t\le1/d$, the smallest eigenvalue of
\[
    \left(I-t\ketbra{\Omega}{\Omega}\right)^{\otimes2}
\]
is
\[
    (1-dt)^2.
\]
By \cref{eq:q2-vectorization}, the displayed infimum is the smallest
eigenvalue, and therefore
\[
    \inf_{C\ne0}
    \frac{q^{(2)}(-t,C)}{\norm{C}_2^2}
    =(1-dt)^2.
\]
This generally differs from $(1-t)(1-dt)$, the expression obtained by
formally substituting $d$ for $r$ in the first branch of
\cref{eq:exact-q2-score}.
\end{remark}

\begin{corollary}[An explicit adjacent-Schmidt-rank violation]
\label[corollary]{cor:adjacent-rank-gap}
Let $1\le r<d$ and put
\[
    Z_{r,d}
    :=
    \left(
      I-\frac1r\ketbra{\Omega}{\Omega}
    \right)^{\otimes2}.
\]
Across the regrouped bipartition $(A_1A_2):(B_1B_2)$,
\[
    \ip{\psi}{Z_{r,d}\psi}\ge0
    \qquad
    \text{whenever }\norm{\psi}=1
    \text{ and }\SR(\psi)\le r.
\]
On the other hand, there is a unit vector $\psi_{r+1}$ of Schmidt
rank $r+1$ such that
\begin{equation}
\label{eq:adjacent-rank-negative-score}
    \ip{\psi_{r+1}}{Z_{r,d}\psi_{r+1}}
    =-\frac1r.
\end{equation}
\end{corollary}

\begin{proof}
The nonnegativity assertion is
\cref{cor:two-copy-block-positive} at $t=1/r$.
Choose a rank-$(r+1)$ projection $P_{r+1}$ and orthogonal unit vectors
$v,w\in\C^d$, and define
\[
    C=P_{r+1}\otimes\ketbra{v}{w},
    \qquad
    \psi_{r+1}:=\frac{\operatorname{vec}(C)}{\sqrt{r+1}}.
\]
Then $\SR(\psi_{r+1})=\rank C=r+1$, while
\[
    \frac{q^{(2)}(-1/r,C)}{\norm{C}_2^2}
    =1-\frac{r+1}{r}
    =-\frac1r.
\]
Now use \cref{eq:q2-vectorization}.
\end{proof}

\begin{remark}[A strict signed separation]
\label[remark]{rem:strict-rank-gap}
Let $0<\varepsilon<1/(r+1)$ and set
$t=(1-\varepsilon)/r$.  Then
\cref{prop:exact-q2-score} gives
\[
    \frac{q^{(2)}(-t,C)}{\norm{C}_2^2}
    \ge
    \varepsilon\left(1-\frac{1-\varepsilon}{r}\right)>0
    \qquad(\rank C\le r),
\]
whereas the rank-$(r+1)$ matrix used in
\cref{cor:adjacent-rank-gap} has normalized score
\[
    1-(r+1)t
    =\frac{(r+1)\varepsilon-1}{r}<0.
\]
Thus the filter has a uniform positive margin on all normalized
matrices of rank at most $r$, while the displayed rank-$(r+1)$ matrix
has negative score.
\end{remark}

This yields the exact $r$-block-positivity threshold for the
two-copy family $(\rho_\alpha^\Gamma)^{\otimes2}$.

\begin{corollary}[Two-copy block positivity]
\label[corollary]{cor:two-copy-block-positive}
Let $d\ge2$ and $\alpha\in\mathbb R$. For every $1\le r\le d^2$, the operator
\[
    (\rho_\alpha^\Gamma)^{\otimes2}
\]
is $r$-block-positive across the regrouped bipartition
\[
    (A_1A_2):(B_1B_2)
\]
if and only if
\begin{equation}
\label{eq:block-positive-threshold}
    \alpha
    \ge
    -\frac1{\min\{r,d\}}.
\end{equation}
\end{corollary}

\begin{proof}
It remains only to consider $\alpha=-t<0$.

Suppose first that $r\le d$. If $t\le1/r$, then
\cref{eq:q2-final-bound} gives
\[
    q^{(2)}(-t,C)\ge0
\]
for every $C$ of rank at most $r$, proving $r$-block positivity.

For necessity in all cases, set
\[
    s:=\min\{r,d\}.
\]
Choose a rank-$s$ projection $P_s$ and orthogonal unit vectors
$v,w\in\C^d$, and let
\[
    C=P_s\otimes\ketbra{v}{w}.
\]
Then
\[
\begin{aligned}
    \rank C&=s\le r,\\
    \Tr C&=0,\\
    \norm{C}_2^2&=s,\\
    \norm{\Tr_1C}_2^2&=s^2,\\
    \norm{\Tr_2C}_2^2&=0.
\end{aligned}
\]
Hence
\[
    q^{(2)}(-t,C)
    =
    s(1-st),
\]
which is negative whenever $t>1/s$. Thus
$r$-block positivity forces
\[
    t\le\frac1{\min\{r,d\}}.
\]

Finally, suppose $r>d$. If $t\le1/d$, then the Choi operator of
\[
    \mathcal R_{-t}(X)
    :=
    \Tr(X)I-tX
\]
is
\[
    J(\mathcal R_{-t})
    =
    I-t\ketbra{\Omega}{\Omega}
    \succeq0,
\]
because $\norm{\Omega}^2=d$ and the only potentially nonpositive
eigenvalue is $1-td$. Therefore $\mathcal R_{-t}$ is completely
positive, so
$\mathcal R_{-t}^{\otimes2}$ is completely positive and hence
$r$-positive for every $r$. This proves sufficiency when $r>d$.
\end{proof}

\subsection{\texorpdfstring{$r$}{r}-positivity of the tensor square}

The preceding block-positivity statement has an equivalent map
formulation through the Choi--Jamiołkowski correspondence: a map is
$r$-positive if and only if its Choi operator is $r$-block-positive
\cite{Choi1975,TerhalHorodecki2000}. With the output factors grouped
against the input factors, the Choi operator of
$\mathcal R_\alpha^{\otimes2}$ is exactly
$(\rho_\alpha^\Gamma)^{\otimes2}$ across the bipartition
$(A_1A_2):(B_1B_2)$ used above.

\begin{corollary}[Tensor-square stability]
\label[corollary]{cor:tensor-square}
Let $d\ge2$ and $\alpha\in\mathbb R$. Define
\[
    \mathcal R_\alpha(X)
    :=
    \Tr(X)I+\alpha X.
\]
For every $1\le r\le d^2$, the tensor-square map
\[
    \mathcal R_\alpha^{\otimes2}
\]
is $r$-positive if and only if
\[
    \alpha
    \ge
    -\frac1{\min\{r,d\}}.
\]
\end{corollary}

\subsection{Exact asymmetric tensor-product theorem}

For $a,b\ge0$, define
\begin{equation}
\label{eq:q2-asymmetric}
    q_{a,b}(C)
    :=
    \norm{C}_2^2
    -a\norm{\Tr_1C}_2^2
    -b\norm{\Tr_2C}_2^2
    +ab\abs{\Tr C}^2.
\end{equation}
Under the same vectorization and regrouping as above,
\[
    q_{a,b}(C)
    =
    \left\langle
      \operatorname{vec}(C),
      \left(I-a\ketbra{\Omega}{\Omega}\right)
      \otimes
      \left(I-b\ketbra{\Omega}{\Omega}\right)
      \operatorname{vec}(C)
    \right\rangle.
\]

\begin{lemma}[One-sided rank bound]
\label[lemma]{lem:one-sided-partial-trace}
If $\rank C\le r$, then
\begin{equation}
\label{eq:one-sided-partial-trace}
    \norm{\Tr_jC}_2^2
    \le r\norm{C}_2^2,
    \qquad j\in\{1,2\}.
\end{equation}
\end{lemma}

\begin{proof}
By trace-norm duality,
\[
    \norm{\Tr_j C}_1
    =
    \sup_{\norm{Y}_\infty\le1}
    \abs{\Tr\left[Y^*\Tr_j C\right]}
    =
    \sup_{\norm{Y}_\infty\le1}
    \abs{\Tr\left[(I\otimes Y^*)C\right]}
    \le\norm{C}_1,
\]
with the identity factor placed on the traced subsystem.  Hence
\[
    \norm{\Tr_j C}_2
    \le\norm{\Tr_j C}_1
    \le\norm{C}_1
    \le\sqrt{r}\norm{C}_2.
\]
\end{proof}

\begin{theorem}[Exact asymmetric tensor-product score]
\label[theorem]{thm:exact-asymmetric-score}
Let $d\ge2$, $1\le r\le d$, and $a,b\ge0$.  Put
\[
    m:=\max\{a,b\},
    \qquad
    \ell:=\min\{a,b\}.
\]
Then
\begin{equation}
\label{eq:exact-asymmetric-score}
    \inf_{\substack{C\in M_{d^2}(\C)\setminus\{0\}\\
                     \rank C\le r}}
    \frac{q_{a,b}(C)}{\norm{C}_2^2}
    =
    \begin{cases}
      (1-\ell)(1-rm),&m\le1/r,\\[1mm]
      1-rm,&m\ge1/r.
    \end{cases}
\end{equation}
\end{theorem}

\begin{proof}
By interchanging the two tensor factors, assume $a=m$ and $b=\ell$.
For $C\ne0$ of rank at most $r$, write
\[
    x:=\frac{\norm{\Tr_1C}_2^2}{\norm{C}_2^2},
    \qquad
    y:=\frac{\norm{\Tr_2C}_2^2}{\norm{C}_2^2},
    \qquad
    z:=\frac{\abs{\Tr C}^2}{\norm{C}_2^2}.
\]
By \cref{thm:rank_r,lem:one-sided-partial-trace},
\[
    x+y\le r+\frac zr,
    \qquad
    x\le r,
    \qquad
    0\le z\le r.
\]
Consequently,
\begin{align*}
    ax+by
    &=b(x+y)+(a-b)x\\
    &\le b\left(r+\frac zr\right)+(a-b)r
      =ar+\frac br z,
\end{align*}
and hence
\begin{equation}
\label{eq:asymmetric-z-bound}
    \frac{q_{a,b}(C)}{\norm{C}_2^2}
    \ge
    1-ar+b\left(a-\frac1r\right)z.
\end{equation}
If $a\le1/r$, minimize the right-hand side over $0\le z\le r$
by taking $z=r$; this gives
\[
    1-ar+br\left(a-\frac1r\right)
    =(1-b)(1-ar).
\]
Equality is attained by
$C=P_r\otimes\ketbra{v}{v}$.
If $a\ge1/r$, minimize by taking $z=0$; this gives $1-ar$, with
equality attained by
$C=P_r\otimes\ketbra{v}{w}$ for $v\perp w$.
If $b>a$, interchange the two tensor factors in these constructions.
\end{proof}

\begin{corollary}[Asymmetric tensor-product $r$-positivity threshold]
\label[corollary]{cor:asymmetric-block-positive}
Let $d\ge2$, $1\le r\le d^2$, and $a,b\ge0$.  Then
\[
    \left(I-a\ketbra{\Omega}{\Omega}\right)
    \otimes
    \left(I-b\ketbra{\Omega}{\Omega}\right)
\]
is $r$-block-positive across the regrouped bipartition if and only if
\begin{equation}
\label{eq:asymmetric-threshold}
    \max\{a,b\}
    \le
    \frac1{\min\{r,d\}}.
\end{equation}
Equivalently,
\[
    \mathcal R_{-a}\otimes\mathcal R_{-b}
\]
is $r$-positive if and only if \cref{eq:asymmetric-threshold} holds.
\end{corollary}

\begin{proof}
For $r\le d$, the assertion follows immediately from
\cref{thm:exact-asymmetric-score}.  If $r>d$ and
$a,b\le1/d$, both factors
$I-a\ketbra{\Omega}{\Omega}$ and
$I-b\ketbra{\Omega}{\Omega}$ are positive semidefinite, so their
tensor product is positive semidefinite.  Conversely, if, say,
$a>1/d$, choose
\[
    C=I_d\otimes\ketbra{v}{w},
    \qquad v\perp w.
\]
Then $\rank C=d\le r$ and
$q_{a,b}(C)/\norm{C}_2^2=1-ad<0$.  The case $b>1/d$ follows by
interchanging the tensor factors.  The map statement follows from the
Choi--Jamio\l kowski correspondence.
\end{proof}

\section{Schmidt-Number Criteria and Verification}

\subsection{Quantitative trace-norm separation from bounded Schmidt number}

\begin{corollary}[Trace-norm separation certified by the two-copy witness]
\label[corollary]{cor:two-copy-trace-distance}
Let $1\le r<d$ and let $Z_{r,d}$ be as in
\cref{cor:adjacent-rank-gap}.  Put
\[
    \gamma:=\frac dr-1>0,
    \qquad
    \Delta_{r,d}:=\gamma+\max\{1,\gamma^2\}.
\]
Define the normalized Schmidt-number-$r$ state set
\[
    \mathcal S_r
    :=
    \left\{
      \sigma\succeq0:
      \Tr\sigma=1,\ \SN(\sigma)\le r
    \right\}.
\]
If $\rho$ is a state satisfying
\[
    \Tr(Z_{r,d}\rho)=-\varepsilon<0,
\]
then
\begin{equation}
\label{eq:two-copy-trace-distance}
    \inf_{\sigma\in\mathcal S_r}
    \norm{\rho-\sigma}_1
    \ge
    \frac{2\varepsilon}{\Delta_{r,d}}.
\end{equation}
\end{corollary}

\begin{proof}
Since $Z_{r,d}$ is $r$-block-positive,
\[
    \Tr(Z_{r,d}\sigma)\ge0
    \qquad
    \text{for every }\sigma\in\mathcal S_r.
\]
The operator
$I-r^{-1}\ketbra{\Omega}{\Omega}$ has eigenvalues $1$ and
$1-d/r=-\gamma$.  Hence $Z_{r,d}$ has eigenvalues
\[
    1,
    \qquad
    -\gamma,
    \qquad
    \gamma^2.
\]
Thus its spectral diameter is
$\Delta_{r,d}$.

Fix $\sigma\in\mathcal S_r$.  Then
\[
    \Tr\bigl[Z_{r,d}(\sigma-\rho)\bigr]\ge\varepsilon.
\]
Because $\Tr(\sigma-\rho)=0$, we may subtract from $Z_{r,d}$ the
midpoint of its extreme eigenvalues.  The resulting centered operator
has operator norm $\Delta_{r,d}/2$.  H\"older duality therefore gives
\[
    \varepsilon
    \le
    \frac{\Delta_{r,d}}2\norm{\sigma-\rho}_1.
\]
Taking the infimum over $\sigma\in\mathcal S_r$ proves the claim.
\end{proof}

\subsection{Matrix-valued Schmidt-number criteria}

\begin{corollary}[A positive-map semidefinite cut]
\label[corollary]{cor:positive-map-cut}
Let $1\le r\le d^2$, let $X$ be a finite-dimensional Hilbert space,
and put
\[
    s:=\min\{r,d\},
    \qquad
    \Lambda_r
    :=
    \mathcal R_{-1/s}^{\otimes2}.
\]
If a positive semidefinite operator
$\sigma\in\cL\bigl(X\otimes(\C^d)^{\otimes2}\bigr)$ has Schmidt
number at most $r$ across
$X:(\C^d)^{\otimes2}$, then
\begin{equation}
\label{eq:positive-map-cut}
    (\id_X\otimes\Lambda_r)(\sigma)\succeq0.
\end{equation}
More generally, \cref{eq:positive-map-cut} holds with
$\Lambda_r$ replaced by
$\mathcal R_{-a}\otimes\mathcal R_{-b}$ whenever
$a,b\ge0$ and
\[
    \max\{a,b\}\le\frac1{\min\{r,d\}}.
\]
\end{corollary}

\begin{proof}
By \cref{cor:asymmetric-block-positive}, the indicated map is
$r$-positive.  Decompose $\sigma$ into positive multiples of pure
states of Schmidt rank at most $r$.  For each such vector, its Schmidt
support on $X$ has dimension at most $r$, so positivity follows from
$r$-positivity after restricting the identity ancilla to that support.
Conjugating back into $X$ and summing the resulting positive operators
proves the assertion.
\end{proof}

\begin{remark}[Use as an SDP cut]
Whenever a semidefinite relaxation contains an explicit candidate
state block $\sigma$, \cref{eq:positive-map-cut} is a linear matrix
inequality in the entries of that block.  If those entries are affine
functions of a moment matrix, applying $\id\otimes\Lambda_r$ produces
an affine matrix pencil and therefore a valid additional semidefinite
constraint under a Schmidt-number-$r$ promise.  This is a necessary
condition for that promise, not in general a complete semidefinite
characterization of Schmidt number at most $r$.
\end{remark}

\subsection{Local aggregation of rank witnesses}

\begin{proposition}[Local-channel aggregation]
\label[proposition]{prop:local-witness-aggregation}
Let $\mathcal H_A$ and $\mathcal H_B$ be finite-dimensional Hilbert
spaces.  For each $e$ in a finite set $E$, let
\[
    \Phi_e^A:\cL(\mathcal H_A)\to\cL(\mathcal K_{A,e}),
    \qquad
    \Phi_e^B:\cL(\mathcal H_B)\to\cL(\mathcal K_{B,e})
\]
be quantum channels, let $W_e=W_e^*\in\cL(\mathcal K_{A,e}\otimes
\mathcal K_{B,e})$ be $r$-block-positive, and let
$\lambda_e\ge0$.  Then
\begin{equation}
\label{eq:aggregated-witness}
    H
    :=
    \sum_{e\in E}
    \lambda_e
    (\Phi_e^A\otimes\Phi_e^B)^*(W_e)
\end{equation}
is $r$-block-positive across $\mathcal H_A:\mathcal H_B$.  In
particular,
\[
    \Tr(H\rho)\ge0
    \qquad
    \text{whenever }\SN(\rho)\le r.
\]
\end{proposition}

\begin{proof}
Local quantum channels do not increase Schmidt number: apply their
local Kraus operators to a decomposition into pure states of Schmidt
rank at most $r$. Hence
$\SN((\Phi_e^A\otimes\Phi_e^B)(\rho))\le r$ for every $e$.
Duality and $r$-block positivity then give
\[
    \Tr(H\rho)
    =\sum_{e\in E}\lambda_e
      \Tr\left[W_e(\Phi_e^A\otimes\Phi_e^B)(\rho)\right]
    \ge0.
\]
\end{proof}

\begin{corollary}[A local-Hamiltonian Schmidt-number certificate]
\label[corollary]{cor:local-hamiltonian-rank-certificate}
In \cref{prop:local-witness-aggregation}, suppose each local output
contains two matched $d_e$-dimensional pairs, with $r\le d_e$, and take
\[
    W_e
    =
    \left(
      I-\frac1r\ketbra{\Omega_e}{\Omega_e}
    \right)^{\otimes2}.
\]
Then the operator $H$ in \cref{eq:aggregated-witness} satisfies
\[
    \Tr(H\rho)<0
    \quad\Longrightarrow\quad
    \SN(\rho)\ge r+1.
\]
For literal subsystem tests, the channels $\Phi_e^A$ and $\Phi_e^B$
may be taken to be the corresponding partial traces, so $H$ is a sum
of local terms embedded into the global Hilbert space.  More generally,
any geometric locality of $H$ is inherited from the subsystem supports
of the chosen channels; arbitrary channels need not produce local
Hamiltonian terms.
\end{corollary}

\begin{remark}[Positive-semidefinite local shifts]
Writing $\Psi_e=\Phi_e^A\otimes\Phi_e^B$, the shifted output term
\[
    h_e^{\mathrm{out}}
    :=W_e-\lambda_{\min}(W_e)I_e\succeq0
\]
has a positive pullback $\Psi_e^*(h_e^{\mathrm{out}})$. Since
$\Psi_e^*(I_e)=I$, every normalized state of Schmidt number at most
$r$ satisfies
\[
    \Tr\!\left[
      \sum_e\lambda_e\Psi_e^*(h_e^{\mathrm{out}})\rho
    \right]
    \ge-\sum_e\lambda_e\lambda_{\min}(W_e).
\]
For the filter witnesses above,
$\lambda_{\min}(W_e)=1-d_e/r$. This supplies a known energy baseline
but, by itself, is not an NLTS statement.
\end{remark}


\subsection{A four-qudit Schmidt-number witness}

The signed adjacent-rank gap can be converted into a literal two-outcome
measurement with perfect completeness.  The clean normalization occurs
when the local dimension is one larger than the forbidden Schmidt rank.

Fix $r\ge1$, put $d=r+1$, and write
\[
    Q_j:=\ketbra{\Omega_d}{\Omega_d}_{A_jB_j},
    \qquad
    \Phi_j:=\frac1dQ_j,
    \qquad j\in\{1,2\}.
\]
Thus $\Phi_j$ is the normalized maximally entangled projection on the
$j$th matched pair.  Define
\begin{equation}
\label{eq:pcp-signed-filter}
    Z_r
    :=
    \left(I-\frac1rQ_1\right)
    \left(I-\frac1rQ_2\right)
\end{equation}
and
\begin{equation}
\label{eq:pcp-accept-effect}
    M_r
    :=
    \frac{r}{r+1}(I-Z_r)
    =
    \Phi_1+\Phi_2-\frac{r+1}{r}\Phi_1\Phi_2.
\end{equation}

\begin{theorem}[A perfect-completeness adjacent-rank witness]
\label[theorem]{thm:pcp-rank-verifier}
Across the bipartition
\[
    (A_1A_2):(B_1B_2),
\]
the operator $M_r$ is a valid POVM effect and satisfies
\begin{equation}
\label{eq:pcp-exact-soundness}
    \sup_{\substack{\rho\succeq0,\ \Tr\rho=1\\
                     \SN(\rho)\le r}}
    \Tr(M_r\rho)
    =\frac{r}{r+1}.
\end{equation}
There is a unit vector $\psi_*$ of Schmidt rank exactly $r+1$ such
that
\begin{equation}
\label{eq:pcp-perfect-completeness}
    \ip{\psi_*}{M_r\psi_*}=1.
\end{equation}
Consequently this four-qudit test has perfect completeness, exact
soundness $r/(r+1)$ against states of Schmidt number at most $r$, and
promise gap $1/(r+1)$.

Operationally, the test measures the two commuting binary
projections $\{\Phi_j,I-\Phi_j\}$.  It rejects if neither projection
clicks, accepts if exactly one clicks, and, if both click, accepts with
probability $(r-1)/r$.
\end{theorem}

\begin{proof}
The operator $I-r^{-1}Q_j$ has eigenvalue $1$ on
$\ker\Phi_j$ and eigenvalue
\[
    1-\frac dr=-\frac1r
\]
on $\ran\Phi_j$.  Hence $Z_r$ has eigenvalues
\[
    1,
    \qquad
    -\frac1r,
    \qquad
    \frac1{r^2}.
\]
It follows that $M_r=r(r+1)^{-1}(I-Z_r)$ has eigenvalues
\[
    0,
    \qquad
    1,
    \qquad
    \frac{r-1}{r},
\]
so $0\preceq M_r\preceq I$.  These three eigenvalues also give the
stated classical postprocessing of the two EPR-projection outcomes.

By \cref{cor:two-copy-block-positive}, $Z_r$ is $r$-block-positive.
Therefore every state $\rho$ with $\SN(\rho)\le r$ obeys
\[
    \Tr(M_r\rho)
    =\frac{r}{r+1}\left(1-\Tr(Z_r\rho)\right)
    \le\frac{r}{r+1}.
\]
This bound is attained.  Choose a rank-$r$ projection $P_r$ on the
first $d$-dimensional factor and orthogonal unit vectors $v,w$ on the
second factor.  For
\[
    C_r=P_r\otimes\ketbra{v}{w},
    \qquad
    \psi_r=\frac{\operatorname{vec}(C_r)}{\sqrt r},
\]
\cref{prop:exact-q2-score} gives
$\ip{\psi_r}{Z_r\psi_r}=0$, and hence
$\ip{\psi_r}{M_r\psi_r}=r/(r+1)$.

For perfect completeness, take
\[
    C_*=I_d\otimes\ketbra{v}{w},
    \qquad
    \psi_*:=\frac{\operatorname{vec}(C_*)}{\sqrt d},
    \qquad v\perp w.
\]
Then $\SR(\psi_*)=\rank C_*=d=r+1$, and
\cref{cor:adjacent-rank-gap} gives
\[
    \ip{\psi_*}{Z_r\psi_*}=-\frac1r.
\]
Substitution into \cref{eq:pcp-accept-effect} yields
\cref{eq:pcp-perfect-completeness}.  Equivalently, $\psi_*$ is a
maximally entangled state on the first matched pair tensored with a
product vector orthogonal to the maximally entangled vector on the
second pair, so exactly one EPR projection clicks with certainty.
\end{proof}

\begin{corollary}[Robust separation from bounded Schmidt number]
\label[corollary]{cor:pcp-robust-self-check}
If a state $\rho$ satisfies
\[
    \Tr(M_r\rho)
    \ge
    \frac{r}{r+1}+\varepsilon,
\]
then
\[
    \inf_{\sigma:\,\SN(\sigma)\le r}
    \frac12\norm{\rho-\sigma}_1
    \ge\varepsilon.
\]
Thus acceptance above the soundness threshold certifies not only
Schmidt number at least $r+1$, but trace distance at least
$\varepsilon$ from the entire Schmidt-number-$r$ state set.
\end{corollary}

\begin{proof}
For every admissible $\sigma$,
\[
    \Tr\bigl[M_r(\rho-\sigma)\bigr]\ge\varepsilon.
\]
The variational characterization of trace distance, together with
$0\preceq M_r\preceq I$, gives the claim.
\end{proof}

\begin{proposition}[Optimal affine conversion of the signed score]
\label[proposition]{prop:pcp-affine-optimal}
Among all POVM effects of the form
\[
    M=aI-bZ_r,
    \qquad b\ge0,
\]
the largest possible acceptance-probability separation between a
state of score $0$ and the explicit state of score $-1/r$ is
$1/(r+1)$.  The effect $M_r$ in
\cref{eq:pcp-accept-effect} attains this optimum.
\end{proposition}

\begin{proof}
Because the extreme eigenvalues of $Z_r$ are $1$ and $-1/r$, the
constraints $0\preceq aI-bZ_r\preceq I$ imply
\[
    a-b\ge0,
    \qquad
    a+\frac br\le1.
\]
Subtracting gives $b(1+1/r)\le1$, hence the separation between scores
$0$ and $-1/r$ is at most $b/r\le1/(r+1)$.  For $M_r$ one has
$a=b=r/(r+1)$, so equality holds.
\end{proof}

\begin{remark}[Interpretation as a verifier gadget]
As a witness, the test unconditionally certifies Schmidt number at
least $r+1$ whenever its threshold $r/(r+1)$ is exceeded. It accepts
one explicit Schmidt-rank-$(r+1)$ family perfectly, but is not a
complete test for all states of that rank. Its interpretation as a
promise-gap verifier additionally requires a reduction or promise
restricting soundness strategies to Schmidt number at most $r$; by
itself it is not a PCP theorem.
\end{remark}

\section{Higher-Copy Filters}

Beyond the exact one- and two-copy cases, the unresolved parameter
regime is
\[
    m\ge3,
    \qquad
    1<r<d,
    \qquad
    \frac1d<t\le\frac1r.
\]
At the endpoint $r=2$, $d=3$, and $t=1/2$, this contains the open
higher-copy Werner-state distillability problem. The results below
therefore provide exact reformulations, propagation rules, and solved
complementary regimes; they do not claim a general higher-copy
threshold theorem.

Let $m\ge1$, let
\[
    \mathcal K=(\C^d)^{\otimes m},
\]
and for each matched output--input pair put
\[
    Q_j:=\ketbra{\Omega_d}{\Omega_d}_{A_jB_j}.
\]
For $t\ge0$, define
\begin{equation}
\label{eq:mfold-filter}
    Z_t^{(m)}
    :=
    \bigotimes_{j=1}^m(I-tQ_j).
\end{equation}
If $S\subseteq[m]$, then $\Tr_S C$ denotes the partial trace of
$C\in\cL(\mathcal K)$ over the tensor factors indexed by $S$, with
$\Tr_\varnothing C=C$.

\begin{proposition}[Inclusion--exclusion and multiplicativity]
\label[proposition]{prop:mfold-inclusion-exclusion}
For every $C\in\cL(\mathcal K)$,
\begin{equation}
\label{eq:mfold-inclusion-exclusion}
\begin{aligned}
    q_m(t,C)
    &:=%
    \ip{\operatorname{vec}(C)}
       {Z_t^{(m)}\operatorname{vec}(C)}\\
    &=
    \sum_{S\subseteq[m]}
    (-t)^{\abs S}\norm{\Tr_S C}_2^2.
\end{aligned}
\end{equation}
Moreover, for operators $C$ and $D$ on disjoint collections of
copies,
\begin{equation}
\label{eq:mfold-multiplicativity}
    q_{m+n}(t,C\otimes D)
    =q_m(t,C)q_n(t,D).
\end{equation}
\end{proposition}

\begin{proof}
For every subset $S$, repeated use of the one-pair vectorization
identity gives
\[
    \left\langle
      \operatorname{vec}(C),
      \left(\prod_{j\in S}Q_j\right)
      \operatorname{vec}(C)
    \right\rangle
    =\norm{\Tr_S C}_2^2.
\]
Expanding the tensor product in \cref{eq:mfold-filter} proves
\cref{eq:mfold-inclusion-exclusion}.  Equation
\cref{eq:mfold-multiplicativity} follows either from this subset
formula or directly from
$Z_t^{(m+n)}=Z_t^{(m)}\otimes Z_t^{(n)}$.
\end{proof}

The scalar expansion has a useful matrix-valued strengthening.  Put
\[
    \mathcal W_r(X):=\Tr(X)I-\frac1rX,
    \qquad
    \Theta_r:=\mathcal W_r\circ T,
    \qquad
    \Theta_r(X)=\Tr(X)I-\frac1rX^T.
\]
The Choi operator of $\Theta_r$ is
$I-r^{-1}F\succeq0$, so $\Theta_r$ is completely positive.

Let $Q=\C^r$, let $E:Q\to\mathcal K$ be an isometry with
$Eq_i=e_i$, and define
\[
    \psi_E=\frac1{\sqrt r}\sum_{i=1}^r e_i\otimes q_i,
    \qquad
    \rho_E=\ketbra{\psi_E}{\psi_E}.
\]
Let $P_\pm(E)$ be the symmetric and antisymmetric projections on
$\overline{E(Q)}\otimes Q$, transported by the identification
$q_i\mapsto\overline e_i$, and extended by zero on the orthogonal
complement of $\overline{E(Q)}\otimes Q$. Then
\[
    \rho_E^{T_{\mathcal K}}
    =\frac1r\bigl(P_+(E)-P_-(E)\bigr).
\]

\begin{proposition}[Exact code-state certificate for every tensor power]
\label[proposition]{prop:mfold-code-certificate}
Define
\begin{equation}
\label{eq:mfold-code-operator}
    A_E^{(m)}
    :=
    (\Theta_r^{\otimes m}\otimes\id_Q)
    \bigl(P_+(E)-P_-(E)\bigr).
\end{equation}
Then
\begin{equation}
\label{eq:mfold-code-channel-form}
    A_E^{(m)}
    =
    r(\mathcal W_r^{\otimes m}\otimes\id_Q)(\rho_E).
\end{equation}
Writing
$\rho_{SQ}:=\Tr_{S^c}\rho_E$, one also has
\begin{equation}
\label{eq:mfold-code-subset-form}
    A_E^{(m)}
    =
    r\sum_{S\subseteq[m]}
    \left(-\frac1r\right)^{\abs S}
    \rho_{SQ}\otimes I_{S^c},
\end{equation}
where, after forming each tensor product, the factors are restored to
their original labeled order.

If
\[
    C=\sum_{i=1}^r\ketbra{e_i}{d_i},
    \qquad
    \delta=\sum_{i=1}^r d_i\otimes q_i,
\]
where zero columns are allowed, then
\begin{equation}
\label{eq:mfold-code-contraction}
    q_m(1/r,C)=\ip{\delta}{A_E^{(m)}\delta}.
\end{equation}
Consequently,
$Z_{1/r}^{(m)}$ is $r$-block-positive if and only if
\[
    A_E^{(m)}\succeq0
\]
for every $r$-dimensional code isometry $E$.
\end{proposition}

\begin{proof}
Since $\mathcal W_r=\Theta_r\circ T$,
\[
\begin{aligned}
    (\mathcal W_r^{\otimes m}\otimes\id_Q)(\rho_E)
    &=(\Theta_r^{\otimes m}\otimes\id_Q)
      (\rho_E^{T_{\mathcal K}})\\
    &=\frac1rA_E^{(m)},
\end{aligned}
\]
which proves \cref{eq:mfold-code-channel-form}.  Expanding every
factor of $\mathcal W_r$ proves
\cref{eq:mfold-code-subset-form}.  The same contraction calculation
as in \cref{lem:partial-trace-contractions}, applied to each subset
$S$, gives \cref{eq:mfold-code-contraction}.  Finally, every matrix of
rank at most $r$ admits such a range factorization after zero padding,
and every choice of $E$ and $\delta$ produces a matrix of rank at most
$r$.  This proves the equivalence.
\end{proof}

\begin{proposition}[Tensor propagation]
\label[proposition]{prop:mfold-tensor-propagation}
Split the copies into disjoint sets $A\sqcup B$.  Suppose the code has
the form
\[
    Eq_i=f_i\otimes\omega,
\]
where $q_i\mapsto f_i$ is a code isometry $E_A$ on the $A$ copies and
$\omega$ is a fixed unit vector on the $B$ copies.  Then
\begin{equation}
\label{eq:mfold-tensor-propagation}
    A_E^{(\abs A+\abs B)}
    =
    A_{E_A}^{(\abs A)}
    \otimes
    \mathcal W_r^{\otimes\abs B}
    (\ketbra{\omega}{\omega}).
\end{equation}
The spectator factor is positive semidefinite.  Hence every
lower-copy positive certificate and every lower-copy kernel vector
propagates after adjoining a fixed spectator state.
\end{proposition}

\begin{proof}
The code state factorizes as
$\rho_E=\rho_{E_A}\otimes\ketbra{\omega}{\omega}$.  Apply
\cref{eq:mfold-code-channel-form} and factor the tensor-power map.
Moreover,
\[
    \mathcal W_r^{\otimes\abs B}(\ketbra{\omega}{\omega})
    =
    \Theta_r^{\otimes\abs B}
    (\ketbra{\overline\omega}{\overline\omega})
    \succeq0
\]
because $\Theta_r$ is completely positive.
\end{proof}

\begin{theorem}[All-copy exact regimes and a universal block-positivity radius]
\label[theorem]{thm:mfold-regimes}
Let $m\ge1$, $d\ge2$, $t\ge0$, and $1\le r\le d^m$.

\begin{enumerate}
\item If $r=1$, then $Z_t^{(m)}$ is $1$-block-positive if and only if
      $t\le1$.
\item If $r\ge d$, then $Z_t^{(m)}$ is $r$-block-positive if and only
      if $t\le1/d$.
\item For $m=1$ and $m=2$, the exact threshold is
      \[
          t\le\frac1{\min\{r,d\}}.
      \]
\item For arbitrary $m$ and $r$, define
      \begin{equation}
      \label{eq:odd-subset-polynomial}
          \vartheta_m(t)
          :=
          \sum_{\substack{0\le k\le m\\k\text{ odd}}}
          \binom mk t^k
          =\frac{(1+t)^m-(1-t)^m}{2}.
      \end{equation}
      If
      \begin{equation}
      \label{eq:mfold-safe-radius}
          r\vartheta_m(t)\le1,
      \end{equation}
      then $Z_t^{(m)}$ is $r$-block-positive.
\item For every $m$, $r$-block positivity forces
      \[
          t\le\frac1{\min\{r,d\}}.
      \]
\end{enumerate}
\end{theorem}

\begin{proof}
For the first assertion, partial transpose on all $B$ systems gives
\[
    \left(Z_t^{(m)}\right)^{\Gamma_B}
    =\bigotimes_{j=1}^m(I-tF_j).
\]
This is positive semidefinite exactly when $t\le1$.  If
$\xi\otimes\eta$ is a product vector across the regrouped cut, then
\[
    \ip{\xi\otimes\eta}{Z_t^{(m)}(\xi\otimes\eta)}
    =
    \ip{\xi\otimes\overline\eta}
       {\left(Z_t^{(m)}\right)^{\Gamma_B}
        (\xi\otimes\overline\eta)}
    \ge0.
\]
Conversely, for $t>1$, take a rank-one tensor product matrix with one
factor a rank-one projection and all remaining factors traceless
rank-one operators; \cref{eq:mfold-multiplicativity} gives normalized
score $1-t<0$.

If $r\ge d$ and $t\le1/d$, every factor
$I-tQ_j$ is positive semidefinite, so the full tensor product is
positive semidefinite.  If $t>1/d$, choose a rank-$d$ projection
$P_d$ in one factor and traceless rank-one operators in every other
factor.  The resulting matrix has rank $d\le r$ and normalized score
$1-dt<0$.

The one-copy threshold is the trace--rank inequality.  The two-copy
threshold is \cref{cor:two-copy-block-positive}.

For the universal sufficient condition, if $C$ has rank at most $r$,
then \cref{lem:one-sided-partial-trace}, applied to the bipartition
$S:S^c$, gives
\[
    \norm{\Tr_S C}_2^2\le r\norm C_2^2
    \qquad(S\ne\varnothing).
\]
Dropping all positive even-cardinality terms in
\cref{eq:mfold-inclusion-exclusion} therefore gives
\[
\begin{aligned}
    q_m(t,C)
    &\ge
    \norm C_2^2
    -\sum_{\abs S\text{ odd}}
       t^{\abs S}\norm{\Tr_S C}_2^2\\
    &\ge
    \left(1-r\vartheta_m(t)\right)\norm C_2^2.
\end{aligned}
\]
This proves \cref{eq:mfold-safe-radius}.

Finally, put $s=\min\{r,d\}$.  Choose a rank-$s$ projection $P_s$ in
one factor and traceless rank-one operators $R_j$ of Hilbert--Schmidt
norm one in all other factors.  By multiplicativity,
\[
    \frac{q_m(t,P_s\otimes R_2\otimes\cdots\otimes R_m)}
         {\norm{P_s\otimes R_2\otimes\cdots\otimes R_m}_2^2}
    =1-st.
\]
It is negative for $t>1/s$, proving necessity.
\end{proof}


\begin{theorem}[Exact all-copy threshold for product-basis diagonal witnesses]
\label[theorem]{thm:mfold-diagonal-threshold}
Fix $m\ge1$ and $d\ge2$.  Restrict $C\in\cL((\C^d)^{\otimes m})$
to operators diagonal in a product basis.

If $r\ge2$, then
\begin{equation}
\label{eq:mfold-diagonal-threshold}
    q_m(t,C)\ge0
    \quad\text{for every such $C$ with $\rank C\le r$}
    \quad\Longleftrightarrow\quad
    t\le\frac1{\min\{r,d\}}.
\end{equation}
For $r=1$, every nonzero diagonal rank-one operator has score
$(1-t)^m\norm C_2^2$; hence the diagonal rank-one condition holds
exactly for $0\le t\le1$ when $m$ is odd, and for every $t\ge0$ when
$m$ is even.
\end{theorem}

\begin{proof}
By local unitary conjugation, use the standard product basis indexed by
words $x=(x_1,\ldots,x_m)\in[d]^m$.  Write
\[
    C=\sum_{x\in T}c_x\ketbra{x}{x},
    \qquad \abs T\le r.
\]
Expanding \cref{eq:mfold-inclusion-exclusion}, or evaluating the EPR
projections directly, gives
\begin{equation}
\label{eq:mfold-diagonal-kernel}
    q_m(t,C)=c^*G_T(t)c,
    \qquad
    G_T(t)_{xy}
    =\prod_{j=1}^m(\delta_{x_jy_j}-t).
\end{equation}
Thus $G_T(t)$ is a principal submatrix of
\[
    (I_d-t\mathbf1\mathbf1^*)^{\otimes m}.
\]

Suppose first that $2\le r<d$ and $t\le1/r$.  For every $x\in T$,
\[
    G_T(t)_{xx}=(1-t)^m.
\]
If $x\ne y$ and $h(x,y)$ is their Hamming distance, then
\[
    \abs{G_T(t)_{xy}}
    =t^{h(x,y)}(1-t)^{m-h(x,y)}
    \le t(1-t)^{m-1},
\]
because $t\le1/2$.  Therefore
\[
    \sum_{y\ne x}\abs{G_T(t)_{xy}}
    \le(\abs T-1)t(1-t)^{m-1}
    \le(1-t)^m.
\]
The Hermitian matrix $G_T(t)$ is diagonally dominant with nonnegative
diagonal, hence positive semidefinite.

If $r\ge d$ and $t\le1/d$, the one-coordinate matrix
$I_d-t\mathbf1\mathbf1^*$ is positive semidefinite.  Its tensor power
and every principal submatrix are therefore positive semidefinite.
This proves sufficiency in \cref{eq:mfold-diagonal-threshold}.

For necessity, put $s=\min\{r,d\}\ge2$.  If $1/s<t<1$, take $T$ to
be $s$ words that vary only in the first coordinate.  Then
\[
    G_T(t)=(1-t)^{m-1}
           (I_s-t\mathbf1\mathbf1^*),
\]
whose all-ones eigenvalue is
$(1-t)^{m-1}(1-st)<0$.  If $t\ge1$, take two words that differ in
every coordinate.  Their $2\times2$ principal matrix has diagonal
$(1-t)^m$ and off-diagonal $(-t)^m$.  Since
$t^m>\abs{1-t}^m$ for $t>1$, it has a negative eigenvalue; at $t=1$
it has eigenvalues $\pm1$.  Both witnesses have support size at most
$r$.  This proves necessity.

Finally, when $r=1$, a diagonal rank-one operator is supported on one
word, and \cref{eq:mfold-diagonal-kernel} reduces to the scalar
$(1-t)^m\norm C_2^2$.
\end{proof}

\begin{remark}[Classical versus collective amplitudes]
The theorem closes the endpoint $t=1/r$ for every copy number on the
product-basis diagonal, hence classical-amplitude, subspace.  Any
failure of the general $m$-copy endpoint for $1<r<d$ must therefore
use genuinely off-diagonal coherent amplitudes; it cannot already be
witnessed by a sparse classical table.
\end{remark}

\begin{remark}[A literal rank-one SOS for every copy number]
Let
\[
    P_\pm^{(j)}:=\frac12(I\pm F_j).
\]
For $0\le t\le1$ and a rank-one matrix $C=\ketbra{x}{y}$,
partial transpose and the spectral decomposition of each flip give
\begin{equation}
\label{eq:mfold-rankone-sos}
\begin{aligned}
    q_m(t,C)
    ={}&
    \sum_{\varepsilon\in\{+,-\}^m}
    (1-t)^{n_+(\varepsilon)}
    (1+t)^{n_-(\varepsilon)}\\
    &\qquad\qquad\times
    \norm{
      \left(\bigotimes_{j=1}^mP_{\varepsilon_j}^{(j)}\right)
      (x\otimes y)
    }^2.
\end{aligned}
\end{equation}
This is an ordinary sum-of-squared-norms certificate for the all-copy
rank-one case.
\end{remark}

At the endpoint $t=1/r$, the code operator
\cref{eq:mfold-code-operator} is the exact object that must be shown
positive for every code $E$.

\begin{proposition}[An $m$-pair verifier and the higher-copy endpoint]
\label[proposition]{prop:mfold-pair-verifier}
Fix $r\ge1$, put $d=r+1$, and define
\[
    M_r^{(m)}
    :=
    \frac{r}{r+1}\left(I-Z_{1/r}^{(m)}\right).
\]
Then $M_r^{(m)}$ is a POVM effect.  If $k$ of the $m$ normalized EPR
projections click, its conditional acceptance probability is
\begin{equation}
\label{eq:mfold-conditional-acceptance}
    a_k
    =
    \frac{r}{r+1}
    \left(1-\left(-\frac1r\right)^k\right).
\end{equation}
There is a Schmidt-rank-$(r+1)$ state accepted with probability one.
Moreover, the following are equivalent:
\begin{enumerate}
\item $Z_{1/r}^{(m)}$ is $r$-block-positive;
\item every state of Schmidt number at most $r$ is accepted by
      $M_r^{(m)}$ with probability at most $r/(r+1)$.
\end{enumerate}
Whenever these conditions hold, the soundness value is exactly
$r/(r+1)$, so the test has promise gap $1/(r+1)$.  Thus the same
perfect-completeness test extends to $m$ matched pairs exactly
when the $m$-copy endpoint is $r$-block-positive.
\end{proposition}

\begin{proof}
Each factor of $Z_{1/r}^{(m)}$ has eigenvalues $1$ and $-1/r$.
Hence the spectrum of the tensor product consists of
$(-1/r)^k$, $0\le k\le m$.  Its largest eigenvalue is $1$ and its
smallest eigenvalue is $-1/r$, so
$0\preceq M_r^{(m)}\preceq I$.  Jointly measuring the commuting EPR
projections gives \cref{eq:mfold-conditional-acceptance}.

Take a maximally entangled vector on one matched pair and, on every
other pair, a product vector orthogonal to the maximally entangled
vector.  The resulting state has Schmidt rank $d=r+1$, exactly one
EPR projection clicks, and $a_1=1$.

Finally,
\[
    \Tr(M_r^{(m)}\rho)
    =\frac{r}{r+1}
      \left(1-\Tr(Z_{1/r}^{(m)}\rho)\right),
\]
so the two assertions are equivalent by convexity of the
Schmidt-number-$r$ state set.  To see that the bound is attained,
choose a rank-$r$ projection $P_r$ on one copy and, on every remaining
copy, a Hilbert--Schmidt-unit rank-one operator $R_j$ with
$\Tr R_j=0$.  Then
\[
    C=P_r\otimes R_2\otimes\cdots\otimes R_m
\]
has rank $r$ and, by \cref{eq:mfold-multiplicativity},
$q_m(1/r,C)=0$.  Its normalized vectorization is therefore accepted
with probability $r/(r+1)$.
\end{proof}

\begin{remark}[The unresolved endpoint]
For $r=2$, the operator $Z_{1/2}^{(m)}$ is the $m$-fold partial
transpose of the Werner endpoint $\rho_{-1/2}$, up to normalization.
Its $2$-block positivity is exactly $m$-copy undistillability.  The
case $m=2$ is proved in this manuscript through the result of Fu, Gao,
and Park, while $m\ge3$ is the remaining higher-copy Werner problem
\cite{fu2026solution2copydistillabilitywerner}.  In particular, a
proof of \cref{prop:mfold-pair-verifier} soundness for
$r=2,d=3,m=3$ would resolve the three-copy qutrit endpoint.
\end{remark}

\section{Further Consequences}

\subsection{Kronecker-sum Ky Fan bounds}

By dualizing the partial-trace inequality over rank-$k$ test
operators, \cref{thm:rank_r} yields a singular-value hierarchy for
Kronecker sums.

\begin{corollary}
\label[corollary]{cor:kronecker}
Let $A,B\in M_d(\C)$ be traceless, not necessarily Hermitian,
operators, and define
\[
    M
    :=
    A\otimes I_d+I_d\otimes B.
\]
For every $1\le k\le d^2$,
\begin{equation}
\label{eq:kronecker-bound}
    \sum_{j=1}^k s_j(M)^2
    \le
    \min\left\{
      d,
      k+\max\left\{
        0,1-\frac{2k}{d}
      \right\}
    \right\}
    \left(
      \norm{A}_2^2+\norm{B}_2^2
    \right).
\end{equation}
In particular, if $d\le2k$, then
\[
    \sum_{j=1}^k s_j(M)^2
    \le
    k\left(
      \norm{A}_2^2+\norm{B}_2^2
    \right).
\]
\end{corollary}

\begin{proof}
The Frobenius-dual form of Ky Fan duality gives
\begin{equation}
\label{eq:ky-fan-duality}
    \left(
      \sum_{j=1}^k s_j(M)^2
    \right)^{1/2}
    =
    \max_{\substack{\rank C\le k\\ \norm{C}_2=1}}
    \abs{\langle C,M\rangle_{\mathrm{HS}}}.
\end{equation}
Since $A$ and $B$ are traceless,
\begin{align}
    \langle C,M\rangle_{\mathrm{HS}}
    ={}&
    \left\langle
      \Tr_2C-\frac{\Tr C}{d}I,
      A
    \right\rangle_{\mathrm{HS}}
    \notag\\
    &+
    \left\langle
      \Tr_1C-\frac{\Tr C}{d}I,
      B
    \right\rangle_{\mathrm{HS}}.
\label{eq:kronecker-pairing}
\end{align}
By Cauchy--Schwarz,
\begin{align}
    \abs{\langle C,M\rangle_{\mathrm{HS}}}^2
    \le{}&
    \left(
      \norm{A}_2^2+\norm{B}_2^2
    \right)
    \notag\\
    &\times
    \left(
      \norm{\Tr_1C}_2^2
      +
      \norm{\Tr_2C}_2^2
      -
      \frac2d\abs{\Tr C}^2
    \right).
\label{eq:kronecker-CS}
\end{align}
For a rank-at-most-$k$ operator $C$ with $\norm{C}_2=1$,
\cref{eq:main-bound} gives
\begin{align}
    &
    \norm{\Tr_1C}_2^2
    +
    \norm{\Tr_2C}_2^2
    -
    \frac2d\abs{\Tr C}^2
    \notag\\
    &\qquad
    \le
    k
    +
    \left(
      \frac1k-\frac2d
    \right)\abs{\Tr C}^2.
\label{eq:kronecker-main-use}
\end{align}
If $1/k-2/d\le0$, the final term may be discarded. If
$1/k-2/d>0$, then
\[
    \abs{\Tr C}^2
    \le
    k\norm{C}_2^2
    =
    k.
\]
Thus the right-hand side of
\cref{eq:kronecker-main-use} is at most
\[
    k+\max\left\{
      0,1-\frac{2k}{d}
    \right\}.
\]
Combining this with
\cref{eq:ky-fan-duality,eq:kronecker-CS} gives the second entry in the
minimum in \cref{eq:kronecker-bound}. On the other hand, tracelessness
implies
\[
    \norm{M}_2^2
    =d\left(\norm{A}_2^2+\norm{B}_2^2\right),
\]
which gives the first entry. This proves \cref{eq:kronecker-bound}.
\end{proof}


\subsection{Exact Schmidt number and witnesses}

For $d\ge2$, let
\[
    \omega_-^{(d)}
    :=
    \frac{P_-^{(d)}}{\binom d2}
\]
denote the normalized antisymmetric projector state on
$\C^d\otimes\C^d$.

\begin{corollary}[Exact Schmidt number]
\label[corollary]{cor:exact-schmidt-number}
With respect to the regrouped bipartition
\[
    (A_1A_2):(B_1B_2),
\]
one has, for every $d_1,d_2\ge2$,
\begin{equation}
\label{eq:SN-four}
    \SN\left(
      \omega_-^{(d_1)}\otimes\omega_-^{(d_2)}
    \right)
    =
    4.
\end{equation}
\end{corollary}

\begin{proof}
Each standard antisymmetric basis vector has Schmidt rank $2$, so
$\omega_-^{(d_1)}$ admits a decomposition into pure states of Schmidt
rank $2$. Moreover, the antisymmetric subspace contains no nonzero
product vector. Hence
\[
    \SN(\omega_-^{(d_1)})=2.
\]
Tensoring Schmidt-rank-$2$ decompositions therefore gives
\[
    \SN\left(
      \omega_-^{(d_1)}\otimes\omega_-^{(d_2)}
    \right)
    \le4.
\]

For the reverse inequality, let
\[
    Q
    :=
    P_-^{(d_1)}\otimes P_-^{(d_2)},
\]
where the first projection acts on $A_1B_1$ and the second on $A_2B_2$.
Under inverse vectorization across
$(A_1A_2):(B_1B_2)$, every vector
$\phi\in\ran Q$ has the form
\[
    \phi=\operatorname{vec}(K),
    \qquad
    K\in
    \mathfrak{so}(\C^{d_1})\otimes\mathfrak{so}(\C^{d_2}).
\]
Thus, for a unit vector $\phi\in\ran Q$, its Schmidt coefficients
$s_1\ge s_2\ge s_3\ge\cdots$ across the regrouped bipartition are the
singular values of a Hilbert--Schmidt unit operator $K$ in the
above double-skew space. The case $d_1\ne d_2$ is included by the
zero-extension argument at the beginning of \cref{lem:double-skew}.
By that lemma,
\[
    s_1^2+s_2^2\le\frac12.
\]
Since $s_3\le s_2\le s_1$,
\[
    2s_3^2
    \le
    s_1^2+s_2^2
    \le
    \frac12.
\]
Consequently,
\[
    s_1^2+s_2^2+s_3^2
    \le
    \frac34.
\]
For any orthogonal projection $Q$ and any $k$, projection duality gives
\begin{align*}
    \sup_{\substack{\norm{z}=1\\ \SR(z)\le k}}
    \ip{z}{Qz}
    &=
    \sup_{\substack{\phi\in\ran Q\\ \norm{\phi}=1}}
    \sup_{\substack{\norm{z}=1\\ \SR(z)\le k}}
    \abs{\ip{z}{\phi}}^2\\
    &=
    \sup_{\substack{\phi\in\ran Q\\ \norm{\phi}=1}}
    \sum_{j=1}^k s_j(\phi)^2;
\end{align*}
see \cite{JohnstonKribs2010}. Applying this identity with $k=3$ gives
\[
    \sup_{\substack{\norm{z}=1\\ \SR(z)\le3}}
    \ip{z}{Qz}
    \le
    \frac34.
\]

On the other hand,
\[
    \omega_-^{(d_1)}\otimes\omega_-^{(d_2)}
    =\frac{Q}{\rank Q},
\]
and hence
\[
    \Tr\left[
      Q\left(
        \omega_-^{(d_1)}\otimes\omega_-^{(d_2)}
      \right)
    \right]
    =1.
\]
If the state had Schmidt number at most $3$, it would admit a convex
decomposition into pure states of Schmidt rank at most $3$
\cite{TerhalHorodecki2000}; its expectation against $Q$ would then be at
most $3/4$, a contradiction. Therefore its Schmidt number is at least
$4$, proving \cref{eq:SN-four}.
\end{proof}

\begin{corollary}[Explicit Schmidt-number witnesses]
\label[corollary]{cor:explicit-witnesses}
For $d_1,d_2\ge2$, let
\[
    Q=P_-^{(d_1)}\otimes P_-^{(d_2)}.
\]
Then
\[
    W_2:=\frac12I-Q
\]
satisfies
\[
    \Tr(W_2\sigma)\ge0
    \qquad
    \text{whenever }\SN(\sigma)\le2,
\]
and
\[
    W_3:=\frac34I-Q
\]
satisfies
\[
    \Tr(W_3\sigma)\ge0
    \qquad
    \text{whenever }\SN(\sigma)\le3.
\]
Thus a negative expectation of $W_2$ certifies Schmidt number at
least $3$, while a negative expectation of $W_3$ certifies Schmidt
number at least $4$. For the state in
\cref{cor:exact-schmidt-number},
\[
\begin{aligned}
    \Tr\left[
      W_2\left(
        \omega_-^{(d_1)}\otimes\omega_-^{(d_2)}
      \right)
    \right]
    &=-\frac12,\\
    \Tr\left[
      W_3\left(
        \omega_-^{(d_1)}\otimes\omega_-^{(d_2)}
      \right)
    \right]
    &=-\frac14.
\end{aligned}
\]
\end{corollary}


\begin{thebibliography}{10}

\bibitem{Choi1975}
Man-Duen Choi.
``Completely positive linear maps on complex matrices.''
\newblock Linear Algebra and its Applications, 10:285--290, 1975.
\newblock \url{https://doi.org/10.1016/0024-3795(75)90075-0}

\bibitem{CostaRico}
Pablo Costa Rico.
``New partial trace inequalities and distillability of Werner states.''
\newblock Letters in Mathematical Physics, 115:47, 2025.
\newblock \url{https://doi.org/10.1007/s11005-025-01935-y}

\bibitem{CostaRicoWolf}
Pablo Costa Rico and Michael M. Wolf.
``Partial trace relations beyond normal matrices.''
\newblock arXiv:2507.18278v1, July 24, 2025.
\newblock \url{https://arxiv.org/abs/2507.18278v1}

\bibitem{fu2026solution2copydistillabilitywerner}
Jinshi Fu, Li Gao, and Sang-Jun Park.
``A solution to 2-copy distillability of Werner states.''
\newblock arXiv:2607.21367v1, July 23, 2026.
\newblock \url{https://arxiv.org/abs/2607.21367v1}

\bibitem{Horodecki1998}
Micha\l{} Horodecki, Pawe\l{} Horodecki, and Ryszard Horodecki.
``Mixed-state entanglement and distillation: Is there a `bound'
entanglement in nature?''
\newblock Physical Review Letters, 80:5239--5242, 1998.
\newblock \url{https://doi.org/10.1103/PhysRevLett.80.5239}

\bibitem{JohnstonKribs2010}
Nathaniel Johnston and David W. Kribs.
``A family of norms with applications in quantum information theory.''
\newblock Journal of Mathematical Physics, 51:082202, 2010.
\newblock \url{https://doi.org/10.1063/1.3459068}

\bibitem{TerhalHorodecki2000}
Barbara M. Terhal and Pawe\l{} Horodecki.
``Schmidt number for density matrices.''
\newblock Physical Review A, 61:040301, 2000.
\newblock \url{https://doi.org/10.1103/PhysRevA.61.040301}

\bibitem{Werner1989}
Reinhard F. Werner.
``Quantum states with Einstein--Podolsky--Rosen correlations admitting a
hidden-variable model.''
\newblock Physical Review A, 40:4277--4281, 1989.
\newblock \url{https://doi.org/10.1103/PhysRevA.40.4277}

\end{thebibliography}

\end{document}
